\documentclass[11pt]{article}

\usepackage[T1]{fontenc}
\usepackage[margin=0.85in]{geometry}
\usepackage{amsmath,amssymb}
\usepackage{booktabs}
\usepackage{longtable}
\usepackage{tabularx}
\usepackage{array}
\usepackage{multirow}
\usepackage{pdflscape}
\usepackage{xcolor}
\usepackage{graphicx}
\usepackage{seqsplit}
\usepackage{microtype}
\usepackage{hyperref}

\hypersetup{
  colorlinks=true,
  linkcolor=blue!55!black,
  urlcolor=blue!55!black,
  citecolor=blue!55!black,
  pdftitle={Reviewer-Directed Rebuttal Experiments for It Cancels},
  pdfauthor={Anonymous Authors}
}

\newcommand{\ours}{\textsc{Ours}}
\newcommand{\refact}{\textsc{Refact}}
\newcommand{\Aw}{A_{\mathrm w}}
\newcommand{\tref}{\tau_{\mathrm{ref}}}
\newcommand{\R}{\mathbb{R}}
\newcommand{\reviewer}[1]{\textbf{Reviewer response:} #1}
\newcommand{\paperchange}[1]{\textbf{Recommended paper change:} #1}
\newcommand{\hexid}[1]{\texttt{\seqsplit{#1}}}
\newcommand{\caveat}[1]{\par\noindent\colorbox{yellow!18}{%
  \parbox{\dimexpr\linewidth-2\fboxsep\relax}{\textbf{Scope caveat.} #1}}\par}

\title{\textbf{Reviewer-Directed Rebuttal Experiments}\\
\large for \emph{It Cancels: \(O(r^2)\) Cholesky Updates for Whitened Operators}\\
\normalsize NeurIPS 2026 Submission \#29421}
\author{Anonymous Authors}
\date{\today}

\begin{document}
\maketitle

\begin{abstract}
We report a reviewer-directed evaluation of the proposed exact
\(O(r^2)\) update for the double-sided whitened operator
\(\Aw=L^{-1}ML^{-\top}\) under rank-one changes to the underlying
positive-definite metric.  The original planned evaluation comprises 145
completed configurations: 48 natural-gradient runs, 86 CPU stability and
representation experiments, and 11 CUDA benchmarks.  The experiments answer
questions about chart-refresh sensitivity, the interpretation of predictive
accuracy, regimes in which one-sided propagation is preferable, stability
under extreme conditioning and near-boundary downdates, CUDA scaling, and the
practical cost of rank-\(k\) extensions.
A post-audit CUDA optimization follow-up, reported separately in E5, replaces
the original large-rank diagnostic with a baseline-compatible parallel
fallback and evaluates a corrected experimental inverse-factor
representation over ten held-out problem seeds, including timing,
repeated-update stability, refresh-cost, and large-rank comparisons.
A separate post-audit E2 rerun uses the recovered UCI Protein data, calibrates
its time ceiling solely on excluded pilot seeds, selects the last completed
training batch without test feedback, and evaluates each fixed seed once.
A provenance-locked E5 extension additionally evaluates seven ranks through
\(65{,}536\) on fixed transitions from The Well's
\texttt{post\_neutron\_star\_merger} trajectory; it was frozen before feature
generation or timing and is reported separately from the synthetic ladder.

The results sharpen rather than uniformly strengthen the original claims.
The refresh threshold controls a clear refresh-frequency--accuracy trade-off.
At matched \(r=1024\), the corrected E2 rerun gives a paired
\ours{}--\refact{} test-MSE difference of
\(0.0051\pm0.0058\), providing no accuracy-win evidence for the exact update.
A thin factor that remains factorized favors one-sided propagation, whereas
explicit dense materialization favors the proposed two-sided representation.
Positive updates remain backward stable but forward error reflects the
conditioning of whitening; downdate amplification exactly follows the
\((1-\ell)^{-1}\) bound.  The fused CUDA kernel benefits strongly from
independent-update batching, and an exact prefix-scan fallback is
\(1.8\)--\(2.1\times\) faster than CUDA refactorization at
\(r=256\)--1024.  A separately labeled inverse-factor representation is
\(3.3\)--\(12.2\times\) faster in held-out cached-state kernel timing at
\(r=256\)--1024.  Its synthetic self-contained four-state ladder reaches
\(71.3\pm0.1\times\) at \(r=65{,}536\), and the independently locked
real-trajectory ladder reaches \(71.755\pm0.225\times\) at the same rank
while passing every predeclared correctness gate.  Repeated-update and
refresh tests show low drift for well-conditioned problems but strong
condition dependence.  Repeated rank-one updates are advantageous for small
\(k\), but their measured crossover occurs much earlier than a FLOP-only
model predicts.  We state below precisely which claims these data support and
which manuscript statements should be weakened.
\end{abstract}

\caveat{The E1 natural-gradient sweep uses a reconstructed, corrected
implementation of the method described in Section~3.5.  Its original copied
artifact lacks the Protein \texttt{.mat} file and exact submitted training
driver, and its stale prototype does not reproduce submitted Setup~B.  E2 is
therefore rerun separately with the recovered Protein file and a
leakage-controlled driver, but that driver is still a new reconstruction
rather than the missing submitted executable.  E1 and E2 are new controlled
evidence, not exact reproductions of the submitted MSE table.}

\section{Scope and reviewer-question map}
\label{sec:scope}

The Area Chair and all three reviewers considered the algebraic contribution
technically sound.  The experiments therefore target the reviewers' stated
scope, interpretation, and empirical-support questions rather than attempting
to replace the submitted evaluation.  Table~\ref{tab:review-map} records the
purpose of each experiment.

\begin{table}[ht]
\centering
\small
\caption{Mapping from experiments to reviewer questions and manuscript
consequences.}
\label{tab:review-map}
\begin{tabularx}{\linewidth}{@{}
  >{\raggedright\arraybackslash}p{0.07\linewidth}
  >{\raggedright\arraybackslash}p{0.22\linewidth}
  >{\raggedright\arraybackslash}p{0.23\linewidth}
  >{\raggedright\arraybackslash}X@{}}
\toprule
ID & Primary question & Reviewer / meta-review & Intended manuscript use \\
\midrule
E1 & Sensitivity to the chart-refresh threshold \(\tref\) &
fdru Q1; Area Chair & Add a threshold-sensitivity table and explain that
\(\tref\) is an application hyperparameter, not an assumption of Theorem~1. \\
E2 & Whether the update itself improves test MSE &
RVfB Q2; Area Chair & Add the leakage-controlled Protein rerun with an
independently calibrated time ceiling; establish fixed-rank equivalence and
withdraw the unsupported accuracy-win interpretation. \\
E3 & When double-sided maintenance is preferable to a one-sided low-rank
factorization & RVfB Q1; relevant to URn7 Q3/Q5 & Add an explicit regime table:
factorized thin consumers favor one-sided propagation; dense or explicitly
materialized consumers favor two-sided maintenance. \\
E4 & Numerical stability near degeneracy and the downdate boundary &
URn7 Q1 & Add forward- and backward-error stress tests, state the role of
\(\|\Aw\|_2\), and recommend reset/guard policies. \\
E5 & GPU scaling and the sequential nature of Givens rotations &
URn7 Q2 & Report supported fused-kernel correctness and batching together
with the tuned, baseline-compatible parallel fallback; report corrected
held-out inverse-factor timing, repeated-update and refresh evidence, and
synthetic plus provenance-locked real-data large-rank ladders as a separately
labeled experimental representation. \\
E6 & Constants and practical cost for rank-\(k\) updates &
URn7 Q4 & Move the full \(O(r^2k+rk^2+k^3)\) cost into the main text and report
measured rather than FLOP-only crossover points. \\
\bottomrule
\end{tabularx}
\end{table}

The experiments do not directly answer fdru's notation and release questions
or URn7's inverse-free and non-whitened questions.  Those should be handled by
textual revisions: define \(G,M,z,w\) before Algorithm~1; state
\(G,M,\Aw\in\R^{r\times r}\), with \(G\) symmetric positive definite and \(M\)
arbitrary; explain one-time whitening for non-whitened models; expand the
inverse-free regime comparison; and give a concrete code-release plan.

\section{Experimental protocol and provenance}
\label{sec:protocol}

\paragraph{Completion.}
All tasks in the original planned matrix completed on their first worker
attempt:
48/48 natural-gradient configurations, 86/86 CPU synthetic configurations,
and 11/11 CUDA configurations.  The CPU and GPU worker error logs are empty.
Preflight smoke tests, which use reduced data or reduced rank and overlap
planned configurations, are not included in inferential tables.
The post-audit E5 campaign also completed: an audit found all 380 tagged JSON
files, exactly ten seeds for every claimed row, and no missing or duplicate
claim cells.  The corrected E2 campaign likewise completed with exactly three
excluded pilot files, one frozen-budget file, twelve evaluation files, and
one summary; every aggregate was independently reconstructed from the raw
JSON.
The locked black-hole extension completed all seven ranks with one result and
one attestation per rank; a strict write-once audit accepted all 70
predeclared real transitions.  One array element had a Slurm
environment-retrieval restart before its batch script or scientific code ran;
that recovery is documented in the provenance appendix.

\paragraph{Software and compute environment.}
All planned runs used PyTorch 2.12.0 with CUDA 12.6 on a Linux compute node.
CUDA experiments ran on an NVIDIA H100 80\,GB HBM3 GPU.  The CPU worker
allocated 100 CPU cores and 128\,GB memory and launched up to ten concurrent
subprocesses without per-run core isolation.  Its queue interleaved the E1
natural-gradient tasks with the single-thread synthetic tasks, so an E1 run
shared the allocation with a varying mixture of other worker processes.  Each
E1 process set PyTorch intra-operation threads to 10 and inter-operation
threads to 1, with \texttt{OMP\_NUM\_THREADS},
\texttt{MKL\_NUM\_THREADS}, and \texttt{OPENBLAS\_NUM\_THREADS} all set to 10.
All 48 planned E1 JSON files record 10 threads, host \texttt{n22}, and Slurm
job \texttt{391303}.  The GPU worker allocated eight H100 GPUs and assigned
one configuration to each GPU.  The post-audit E5 campaign instead ran as
Slurm job \texttt{393591}, a single 33:48 preempt allocation on host
\texttt{n01} with one H100, four CPUs, and 64\,GB host memory
(\(0.5633\) allocated H100-hours).  It executed seed blocks 1--10 serially and
started a fresh Python process for every benchmark-driver invocation; the
refresh driver emits five policy rows per invocation, and each large-rank
paired row combines two separately invoked method processes.  Seed 0 was
reserved for configuration selection.  Corrected E2 job \texttt{393642} ran
on host \texttt{n02}; every process used ten PyTorch threads.  Although its
batch script requested three exclusive seed steps at once, Slurm serialized
all fifteen training steps.  The prior E2 diagnostic job was serialized in
the same way, so concurrency is comparable across those two campaigns but is
one process, not the three claimed by the script comment.
The black-hole ladder used Slurm array \texttt{394033} on seven distinct H100
UUIDs, one per rank, under the \texttt{marlowe-m000151-pm06} allocation.
Every element received 12 CPUs and recorded the same PyTorch
2.12.0+\texttt{cu126}, CUDA 12.6, cuDNN 91002, and Python 3.10.12 stack.

\paragraph{Core correctness checks.}
Every planned benchmark constructs the exact updated metric and operator and
compares the incremental result with fresh Cholesky factorization and
triangular solves.  CPU rank-\(k\) tests use FP64 and require relative errors
below \(10^{-9}\).  Planned CUDA tests use FP32 and require relative errors
below \(2\times10^{-4}\); all planned configurations passed.  The post-audit
inverse-factor follow-ups use the validation scopes stated in
Section~\ref{sec:e5-inverse}: numerical comparison with fresh results through
\(r=32{,}768\), and full finite-output checks at larger measured ranks.
The locked real-data ladder in Section~\ref{sec:e5-blackhole} additionally
retains sixteen fresh cross-method action probes at every rank; all seven rows
passed their predeclared \(10^{-3}\) gate.

\paragraph{Summary statistics.}
Natural-gradient and the original controlled stability sweeps use three
independent seeds.  Corrected E2 uses three excluded calibration seeds
(1000--1002) and three fixed evaluation seeds (0--2); no pilot evaluates the
test set, and each evaluation does so once after time-based checkpoint
selection.  The post-audit E5 campaign uses ten held-out problem seeds
(1--10).  Reported ``\(\pm\)'' values are sample standard deviations.  In
E5's main and cached-state timing tables, each seed-level latency is a
CUDA-event median over 60 prepared repeats; in the large-rank table it is a
median over 10, 5, or 3 prepared repeats as rank increases.  E5 speedups are
formed as paired within-seed ratios before their mean and sample standard
deviation are computed.  The refresh table instead aggregates one timed
200-update sequence per seed and refresh-policy row, plus the shared
fresh-factorization baseline for each seed and rank.  Because job 393591 ran
serially on one node and one GPU, these E5 standard deviations quantify
seed/process/time variation on that device, not device-to-device uncertainty.
The real-data ladder instead treats ten fixed, temporally related transitions
as its \(n=10\) observations: each contributes its within-transition timing
median, paired ratios are formed within transition, and the reported 95\%
intervals use Student's \(t_9\).  Repeat-level timings are not treated as
independent observations.  These intervals describe variation along one
trajectory rather than uncertainty across independent simulations.
Other microbenchmarks report median latency after warm-up; their mutable input
state is prepared outside the timed interval.  The representation and
rank-\(k\) microbenchmarks pin each operation to one CPU thread.

\section{E1: sensitivity to the refresh threshold}
\label{sec:e1}

\reviewer{fdru Q1 asks how gracefully the method responds across chart-refresh
thresholds.  The Area Chair also identifies this question as a priority.}

\subsection{Motivation and method}

The threshold \(\tref\) triggers chart refresh when the fraction of gradient
energy outside the current rank-\(r\) subspace exceeds \(\tref\).  It belongs
to the transported-subspace natural-gradient application, not to the
double-sided update identity itself.  We used the UCI CASP/Protein regression
data, an 80/20 split, ten epochs, batch size 256, and three seeds.  The
preconditioned first layer has 5,120 parameters.  We evaluated
\[
  r\in\{256,1024\},\qquad
  \tref\in\{0.01,0.02,0.05,0.1,0.2,0.5\}.
\]
The corrected optimizer starts from \(G_0=\lambda I\), \({\Aw}_0=0\), uses
\(\lambda=1\), and periodically refactorizes every 50 steps.  At
\(\tref=0.1\), we additionally ran \ours{} at \(r=512\) and \refact{} at all
three ranks.  Each row in Table~\ref{tab:e1} averages three seeds.

\begin{landscape}
\begin{longtable}{@{}llrrrrrr@{}}
\caption{Refresh-threshold sensitivity.  Refresh and refresh-time columns are
fractions of training steps and application wall time, respectively.
``Application wall'' is end-to-end training time on the reconstructed,
shared-worker backend; it compares a serial C++ incremental call with
10-thread PyTorch BLAS/LAPACK inside otherwise common application code and is
not an isolated kernel comparison.  See Table~\ref{tab:e6} for isolated
single-thread timing.  Drift is the maximum relative \(\Aw\) discrepancy
observed immediately before a periodic fresh reset.}
\label{tab:e1}\\
\toprule
Method & \(r\) & \(\tref\) & Refresh & Refresh time & Test MSE &
Application wall (s) & Max.\ drift \\
\midrule
\endfirsthead
\toprule
Method & \(r\) & \(\tref\) & Refresh & Refresh time & Test MSE &
Application wall (s) & Max.\ drift \\
\midrule
\endhead
\ours & 256 & 0.01 & \(0.668\pm0.015\) & \(0.461\pm0.012\) &
\(0.5801\pm0.0084\) & \(9.0\pm0.9\) & \(1.08{\times}10^{-6}\) \\
\ours & 256 & 0.02 & \(0.485\pm0.003\) & \(0.387\pm0.018\) &
\(0.5786\pm0.0120\) & \(8.5\pm1.0\) & \(1.07{\times}10^{-6}\) \\
\ours & 256 & 0.05 & \(0.272\pm0.008\) & \(0.269\pm0.007\) &
\(0.5800\pm0.0141\) & \(7.1\pm0.3\) & \(1.39{\times}10^{-6}\) \\
\ours & 256 & 0.10 & \(0.160\pm0.003\) & \(0.168\pm0.011\) &
\(0.5842\pm0.0162\) & \(5.6\pm0.5\) & \(2.00{\times}10^{-6}\) \\
\ours & 256 & 0.20 & \(0.091\pm0.002\) & \(0.114\pm0.003\) &
\(0.5895\pm0.0057\) & \(5.7\pm0.1\) & \(2.54{\times}10^{-6}\) \\
\ours & 256 & 0.50 & \(0.026\pm0.001\) & \(0.034\pm0.001\) &
\(0.5901\pm0.0220\) & \(5.1\pm0.2\) & \(3.81{\times}10^{-6}\) \\
\midrule
\ours & 512 & 0.10 & \(0.155\pm0.005\) & \(0.267\pm0.011\) &
\(0.5837\pm0.0076\) & \(10.8\pm0.7\) & \(2.51{\times}10^{-6}\) \\
\midrule
\ours & 1024 & 0.01 & \(0.445\pm0.007\) & \(0.670\pm0.006\) &
\(0.5769\pm0.0100\) & \(44.4\pm1.7\) & \(1.87{\times}10^{-6}\) \\
\ours & 1024 & 0.02 & \(0.344\pm0.003\) & \(0.591\pm0.005\) &
\(0.5722\pm0.0105\) & \(40.0\pm1.4\) & \(3.15{\times}10^{-6}\) \\
\ours & 1024 & 0.05 & \(0.223\pm0.009\) & \(0.460\pm0.005\) &
\(0.5768\pm0.0126\) & \(36.3\pm0.1\) & \(3.60{\times}10^{-6}\) \\
\ours & 1024 & 0.10 & \(0.141\pm0.005\) & \(0.350\pm0.013\) &
\(0.5819\pm0.0084\) & \(27.6\pm4.0\) & \(3.17{\times}10^{-6}\) \\
\ours & 1024 & 0.20 & \(0.073\pm0.001\) & \(0.224\pm0.012\) &
\(0.5881\pm0.0059\) & \(22.5\pm0.4\) & \(7.95{\times}10^{-6}\) \\
\ours & 1024 & 0.50 & \(0.018\pm0.001\) & \(0.068\pm0.002\) &
\(0.5964\pm0.0207\) & \(20.0\pm0.7\) & \(1.96{\times}10^{-5}\) \\
\midrule
\refact & 256 & 0.10 & \(0.165\pm0.007\) & \(0.166\pm0.007\) &
\(0.5820\pm0.0091\) & \(7.0\pm0.3\) & --- \\
\refact & 512 & 0.10 & \(0.158\pm0.003\) & \(0.293\pm0.002\) &
\(0.5791\pm0.0137\) & \(9.5\pm0.1\) & --- \\
\refact & 1024 & 0.10 & \(0.141\pm0.002\) & \(0.383\pm0.005\) &
\(0.5870\pm0.0088\) & \(23.2\pm1.5\) & --- \\
\bottomrule
\end{longtable}
\end{landscape}

\subsection{Results and threshold implications}

Thresholding behaves smoothly.  At \(r=256\), raising \(\tref\) from 0.01 to
0.5 reduces the refresh frequency from 66.8\% to 2.6\% and refresh-time share
from 46.1\% to 3.4\%.  At \(r=1024\), the same change reduces refresh
frequency from 44.5\% to 1.8\% and refresh-time share from 67.0\% to 6.8\%.
The observed application wall-time cost of aggressive refreshing grows
strongly with rank under this execution configuration.

Test-MSE intervals overlap broadly across \(\tref\in[0.01,0.1]\).  The lowest
means occur at \(\tref=0.02\), but the experiment does not establish a
statistically significant optimum.  Very infrequent refresh
(\(\tref\in\{0.2,0.5\}\)) tends to increase MSE and reset drift.  Based on the
threading-independent refresh counts and the overlapping MSE intervals,
\(\tref\approx0.05\)--0.1 is a reasonable starting range in this
reconstruction: it reduces refresh events substantially without evidence of
an accuracy loss.  The application seconds and refresh-time shares are
system- and concurrency-dependent, so a deployment should calibrate the
wall-time trade-off on its own backend rather than treating these seconds as
portable.

\subsection{Application timing versus isolated kernel timing}

The ``Application wall'' column in Table~\ref{tab:e1} includes the full
training loop, chart projection and refresh, model forward/backward passes,
preconditioning, and parameter updates.  In the non-refresh core, \ours{}
calls the serial C++ single-update path, whereas \refact{} calls PyTorch
Cholesky and triangular-solve routines with ten intra-operation threads.
Together with the shared worker conditions described in
Section~\ref{sec:protocol}, this makes the column an application measurement,
not a kernel benchmark.

The isolated comparison is the \(r=1024,k=1\) row of
Table~\ref{tab:e6}.  With both methods pinned to one CPU thread in FP64, the
median incremental latency is 7.776\,ms and the fresh-factorization latency is
70.548\,ms, a 9.072\(\times\) speedup.  The fresh path constructs \(G^+\) and
\(M^+\), factors \(G^+\), and applies the two matrix triangular solves needed
to whiten \(M^+\).  State cloning for the incremental path is prepared outside
the timed interval.

\subsection{Matched-rank equivalence}

\reviewer{RVfB Q2 asks whether the submitted ``strictly better test MSE''
statement can be caused by the update identity itself.}

At fixed rank and with identical chart logic, an exact incremental update and
fresh refactorization define the same model up to numerical and
execution-order effects.  The \(\tref=0.1\) rows in
Table~\ref{tab:e1} provide the matched-rank check.  The refresh fractions for
\ours{}/\refact{} are 0.160/0.165, 0.155/0.158, and 0.141/0.141 at
\(r=256,512,1024\), respectively.  Their final test MSEs are
\(0.5842\pm0.0162\)/\(0.5820\pm0.0091\),
\(0.5837\pm0.0076\)/\(0.5791\pm0.0137\), and
\(0.5819\pm0.0084\)/\(0.5870\pm0.0088\).  The seed variation overlaps
broadly at every rank, consistent with algebraic equivalence rather than a
method-induced accuracy advantage.

Because this reconstructed backend does not reproduce the submitted rank--MSE
trend, these values should not replace the submitted application table.  They
support only the causal correction that the exact update does not itself
improve predictive accuracy.

\paperchange{Add a compact version of Table~\ref{tab:e1}, distinguish
\(\tref\) from the periodic numerical reset interval, and describe
rank-aware, system-specific calibration.  Do not claim a universal optimal
threshold.  Remove ``strictly better test MSE'' from the abstract,
introduction, and conclusion; state instead that the proposed update is
algebraically equivalent to fresh refactorization at fixed rank.  Do not add a
budget--accuracy claim from the old checkpoint-truncated analysis; use the
separately calibrated E2 rerun below for the more limited time-ceiling
comparison.}

\section{E2: leakage-controlled Protein time-ceiling rerun}
\label{sec:e2}

\reviewer{RVfB Q2 asks whether the submitted ``strictly better test MSE''
statement is an effect of the exact update itself.  The Area Chair asks that
the accuracy comparison be made auditable.}

\subsection{Frozen protocol and integrity checks}

The corrected rerun uses the recovered UCI Protein file
(SHA-256 prefix \texttt{5ce538fa}), a fixed 80/20 split with split seed 42,
and normalization fitted only on the 36,584 training examples; the 9,146
test examples are transformed using those training statistics.  The model,
batch size 256, ten-epoch maximum, chart threshold \(\tref=0.1\), and
50-step numerical-reset interval are fixed across methods.  The evaluation
grid is \ours{} at \(r=1024\) and \refact{} at
\(r\in\{256,512,1024\}\), with seeds 0--2.

Before any evaluation, three excluded \ours{} \(r=1024\) pilot seeds
1000--1002 run for ten epochs without evaluating test MSE.  Their training
times are 17.840671, 18.562068, and 18.604640\,s, so the frozen ceiling is
their median, 18.562068\,s.  An evaluation retains the last \emph{completed
batch} at or below that ceiling, restores that checkpoint, and then computes
test MSE exactly once.  Static control-flow inspection finds one such
post-selection test expression; the JSON field recording one test evaluation
is metadata rather than an independently instrumented counter.

This protocol removes the old epoch-checkpoint boundary artifact and
within-run test-set checkpoint selection.  Seeds 0--2 and the same test split
had appeared in the prior diagnostic, so this is not a globally untouched
test-set confirmation.  It is nevertheless a time \emph{ceiling}
combined with a ten-epoch cap, not strict equal-compute training: faster
configurations stop after 1,420 steps even when time remains.  The pilot
aggregate is \(18.335793\pm0.429316\)\,s over 1,420 steps, with
\(7.799395\pm0.118849\)\,s in the non-refresh core,
\(197.3\pm5.7\) refreshes taking \(7.025294\pm0.183706\)\,s,
\(0.214082\pm0.003646\)\,s in periodic resets, and maximum reset drift
\((2.499\pm0.355)\times10^{-6}\).

\begin{table}[ht]
\centering
\scriptsize
\caption{Corrected E2 Protein evaluation under an independently calibrated
18.562068\,s training-time ceiling and ten-epoch cap.  Entries are
mean \(\pm\) sample SD over fixed evaluation seeds 0--2.  ``Training'' is
selected training-loop time; ``Core'' is the accumulated non-refresh update
region.  The ceiling makes the application seconds comparable but not an
isolated-kernel benchmark.}
\label{tab:e2-corrected}
{\setlength{\tabcolsep}{3.5pt}
\begin{tabular}{@{}lrrrrrr@{}}
\toprule
Method & \(r\) & Test MSE & Selected steps & Training s & Core s &
Refresh count \\
\midrule
\ours{} & 1024 &
  \(0.564804\pm0.017907\) & \(1415.0\pm8.7\) &
  \(18.4097\pm0.1995\) & \(7.7635\pm0.1882\) &
  \(199.0\pm7.5\) \\
\refact{} & 256 &
  \(0.550845\pm0.015842\) & \(1420.0\pm0.0\) &
  \(4.4786\pm0.1243\) & \(0.6523\pm0.0068\) &
  \(230.7\pm18.9\) \\
\refact{} & 512 &
  \(0.551236\pm0.018148\) & \(1420.0\pm0.0\) &
  \(6.5422\pm0.0966\) & \(1.5968\pm0.0432\) &
  \(219.3\pm16.0\) \\
\refact{} & 1024 &
  \(0.559725\pm0.014636\) & \(1382.3\pm54.4\) &
  \(18.4492\pm0.1863\) & \(8.3169\pm0.2824\) &
  \(200.3\pm11.9\) \\
\bottomrule
\end{tabular}
}
\end{table}

\subsection{Results and interpretation}

At matched \(r=1024\), the paired per-seed
\ours{}--\refact{} test-MSE differences are
\(0.011670\), \(0.002862\), and \(0.000706\), giving
\(0.005079\pm0.005808\).  The corresponding 95\% \(t_2\) interval is
\([-0.00935,0.01951]\).  Thus \refact{} is numerically lower on all three
seeds in this run, but \(n=3\) does not establish an accuracy difference.
The direction also reverses the small, non-significant difference in the
prior diagnostic.  Neither outcome supports an accuracy-win headline.  They
are consistent with the algebraic expectation of fixed-rank equivalence, but
the time ceiling selects different step counts, so E2 alone is not a
controlled causal equivalence test; the matched-epoch E1 rows provide the
cleaner empirical check.

\ours{} descriptively completes \(1415.0\pm8.7\) selected steps at
\(r=1024\), versus
\(1382.3\pm54.4\) for \refact{}, while both consume essentially the frozen
ceiling.  The paired step difference is \(32.7\pm45.8\), so this small
three-seed sample is only descriptive evidence that the optimized
implementation can permit more training work in this reconstructed
ten-thread application, not an isolated-kernel speedup.

\begin{landscape}
{\scriptsize
\setlength{\tabcolsep}{2.4pt}
\begin{longtable}{@{}llrrrrrrrrr@{}}
\caption{Seed-level corrected E2 audit.  Steps and training seconds are
selected/attempted; refresh is count/seconds.  Pilot runs never evaluate the
test set.}
\label{tab:e2-seeds}\\
\toprule
Phase & Method & \(r\) & Seed & Test MSE & Steps & Train s & Core s &
Refresh & Reset s & Max drift \\
\midrule
\endfirsthead
\toprule
Phase & Method & \(r\) & Seed & Test MSE & Steps & Train s & Core s &
Refresh & Reset s & Max drift \\
\midrule
\endhead
pilot & \ours{} & 1024 & 1000 & --- & 1420/1420 & 17.841/17.841 &
  7.682 & 191/6.814 & 0.210 & \(2.77{\times}10^{-6}\) \\
pilot & \ours{} & 1024 & 1001 & --- & 1420/1420 & 18.562/18.562 &
  7.796 & 202/7.142 & 0.217 & \(2.10{\times}10^{-6}\) \\
pilot & \ours{} & 1024 & 1002 & --- & 1420/1420 & 18.605/18.605 &
  7.920 & 199/7.121 & 0.215 & \(2.63{\times}10^{-6}\) \\
\midrule
eval & \ours{} & 1024 & 0 & 0.575657 & 1420/1420 & 18.183/18.183 &
  7.791 & 192/6.814 & 0.210 & \(4.13{\times}10^{-6}\) \\
eval & \ours{} & 1024 & 1 & 0.574619 & 1405/1406 & 18.558/18.567 &
  7.563 & 207/7.405 & 0.215 & \(1.56{\times}10^{-6}\) \\
eval & \ours{} & 1024 & 2 & 0.544136 & 1420/1420 & 18.488/18.488 &
  7.936 & 198/7.066 & 0.212 & \(2.42{\times}10^{-6}\) \\
\midrule
eval & \refact{} & 256 & 0 & 0.538044 & 1420/1420 & 4.537/4.537 &
  0.657 & 216/0.676 & 0 & --- \\
eval & \refact{} & 256 & 1 & 0.568562 & 1420/1420 & 4.563/4.563 &
  0.644 & 252/0.807 & 0 & --- \\
eval & \refact{} & 256 & 2 & 0.545928 & 1420/1420 & 4.336/4.336 &
  0.656 & 224/0.735 & 0 & --- \\
\midrule
eval & \refact{} & 512 & 0 & 0.541386 & 1420/1420 & 6.559/6.559 &
  1.643 & 203/1.774 & 0 & --- \\
eval & \refact{} & 512 & 1 & 0.572179 & 1420/1420 & 6.629/6.629 &
  1.557 & 235/2.020 & 0 & --- \\
eval & \refact{} & 512 & 2 & 0.540144 & 1420/1420 & 6.438/6.438 &
  1.590 & 220/1.884 & 0 & --- \\
\midrule
eval & \refact{} & 1024 & 0 & 0.563987 & 1420/1420 & 18.234/18.234 &
  8.554 & 187/6.519 & 0 & --- \\
eval & \refact{} & 1024 & 1 & 0.571757 & 1320/1321 & 18.561/18.571 &
  8.004 & 210/7.483 & 0 & --- \\
eval & \refact{} & 1024 & 2 & 0.543430 & 1407/1408 & 18.553/18.590 &
  8.393 & 204/7.152 & 0 & --- \\
\bottomrule
\end{longtable}
}
\end{landscape}

The \(r=256\) and \(r=512\) \refact{} runs finish all ten epochs in only 4.48
and 6.54\,s; their unused ceiling illustrates why the result must not be
described as equal-compute training.

The cross-rank \refact{} means are \(0.550845\), \(0.551236\), and
\(0.559725\) at \(r=256,512,1024\): this run does not show monotonic accuracy
improvement with rank.  Paired \(r=1024-r=256\) differences are
\(0.025943\), \(0.003195\), and \(-0.002498\), or
\(0.008880\pm0.015049\), so the apparent ordering is largely one seed and is
not an established rank effect.  The rank-specific subspaces were sampled
independently rather than nested, and learning rate, damping, and threshold
were held fixed rather than tuned per rank.  Mean unexplained gradient energy
also remains in the narrow 6.3--6.7\% range across these runs.  That is
compatible with a low effective dimension, but the thresholded refresh
statistic is not an effective-rank spectrum and cannot establish that
explanation.

Mean refresh time is \(7.09496\pm0.29702\)\,s for \ours{} at \(r=1024\)
and \(0.73923\pm0.06525\), \(1.89269\pm0.12327\), and
\(7.05145\pm0.48989\)\,s for \refact{} at \(r=256,512,1024\),
respectively.  The \ours{} periodic resets add
\(0.21234\pm0.00246\)\,s.

For \ours{}, evaluation reset drift is
\((2.702\pm1.310)\times10^{-6}\), with a maximum of
\(4.131\times10^{-6}\).  No numerical failure occurs.  The optimized
positive-update source path and SHA-256 are recorded in all six \ours{} pilot
and evaluation JSONs; the source hash is
\texttt{b1c92d7f...d87b4}.  The loaded shared-object path is also recorded;
its binary hash was verified externally as \texttt{e500c96a...}, but was not
written into the run JSON and therefore is not self-authenticating
provenance.

\subsection{Implementation diagnostic versus the prior campaign}

Relative to the superseded \texttt{e2proper-dataf073752} campaign, mean
\ours{} \(r=1024\) core time falls from
\(23.1835\pm0.2543\) to \(7.7635\pm0.1882\)\,s, a \(2.986\times\)
reduction, while selected training-loop time falls from
\(33.3943\pm0.3742\) to \(18.4097\pm0.1995\)\,s, a \(1.814\times\)
reduction.  This is diagnostic evidence for the optimized implementation,
not a controlled hardware speedup: the prior campaign ran on host
\texttt{n05}, the corrected campaign ran on \texttt{n02}, and CPU-model
metadata were not recorded.  The prior JSONs also record neither a source nor
binary hash, so the old executable cannot be reconstructed exactly.

The optimized kernel is algebraically equivalent but changes FP32
reassociation.  In the nonlinear thresholded training loop, refresh
decisions and subsequent parameters can therefore diverge.  Prior and current
\ours{} test MSE are \(0.551332\pm0.016693\) and
\(0.564804\pm0.017907\), respectively, including changed outcomes for seeds
that complete all 1,420 steps.  We report this sensitivity rather than
attributing the MSE change to the time ceiling or selecting the favorable
campaign.

\paperchange{Replace the withdrawn checkpoint-truncated E2 table with
Table~\ref{tab:e2-corrected}.  Describe the result as a common ten-epoch
horizon with an independently calibrated upper time cutoff, not strict
equal-compute training.  Remove the accuracy-win claim and state that the
matched-rank paired result is consistent with numerical/execution-order
variation around an exact update.}

\section{E3: one-sided versus two-sided representations}
\label{sec:e3}

\reviewer{RVfB Q1 asks when double-sided maintenance is preferable if
\(M=BC^\top\) is thin and a one-sided factor can be propagated.  The result
also sharpens the regime discussion requested by URn7 Q3 and Q5.}

\subsection{Motivation and method}

We form \(M=BC^\top\) with \(B,C\in\R^{r\times p}\) and compare five
single-thread CPU operations after a rank-one metric update:
\begin{enumerate}
  \item \textbf{Two-sided:} directly update dense \(\Aw\) with the proposed
  \(O(r^2)\) kernel.
  \item \textbf{Drag:} one-sided Seeger-style propagation of
  \(L^{-1}B\) and \(L^{-1}C\), without materializing \(\Aw\).
  \item \textbf{Drag+mat.:} drag both factors and then form their dense
  product.
  \item \textbf{Fresh:} freshly factorize \(G^+\) and triangular-solve for
  both thin factors, without materialization.
  \item \textbf{Fresh+mat.:} freshly solve and then form the dense product.
\end{enumerate}
We test \(r\in\{256,1024\}\), factor widths from \(p=1\) through \(2r\), 30
median-timing repetitions (15 for \(p>r\)), and verify all representations
against the same dense result.  Preparation and state cloning are outside the
timed region.

\begin{landscape}
\scriptsize
\begin{longtable}{@{}rrrrrrrr@{}}
\caption{Representation benchmark.  Times are milliseconds.  ``Factor MiB''
is logical FP32 storage (four bytes per value), counting only the lower
triangle of the Cholesky factor plus both one-sided factors; an implementation
that stores the factor densely may allocate more.  Under the same convention,
the dense two-sided state occupies 0.375\,MiB at \(r=256\) and 6.002\,MiB at
\(r=1024\).}
\label{tab:e3-full}\\
\toprule
\(r\) & \(p\) & Two-sided & Drag & Drag+mat. & Fresh & Fresh+mat. &
Factor MiB \\
\midrule
\endfirsthead
\toprule
\(r\) & \(p\) & Two-sided & Drag & Drag+mat. & Fresh & Fresh+mat. &
Factor MiB \\
\midrule
\endhead
256 & 1   & 0.340 & 0.077 & 0.085 & 0.399 & 0.404 & 0.127 \\
256 & 4   & 0.341 & 0.078 & 0.087 & 0.403 & 0.414 & 0.133 \\
256 & 8   & 0.335 & 0.078 & 0.091 & 0.406 & 0.416 & 0.141 \\
256 & 16  & 0.339 & 0.078 & 0.096 & 0.416 & 0.434 & 0.157 \\
256 & 32  & 0.346 & 0.080 & 0.109 & 0.446 & 0.467 & 0.188 \\
256 & 64  & 0.340 & 0.081 & 0.133 & 0.481 & 0.530 & 0.250 \\
256 & 128 & 0.346 & 0.085 & 0.182 & 0.585 & 0.676 & 0.375 \\
256 & 256 & 0.344 & 0.092 & 0.278 & 0.762 & 0.957 & 0.625 \\
256 & 512 & 0.365 & 0.113 & 0.484 & 1.166 & 1.560 & 1.125 \\
\midrule
1024 & 1    & 6.465 & 2.786 & 2.989 & 15.487 & 17.103 & 2.010 \\
1024 & 4    & 7.911 & 3.049 & 3.523 & 20.651 & 21.113 & 2.033 \\
1024 & 8    & 8.753 & 3.390 & 3.671 & 21.313 & 22.040 & 2.064 \\
1024 & 16   & 9.426 & 3.446 & 4.102 & 22.903 & 23.859 & 2.127 \\
1024 & 32   & 9.135 & 3.451 & 4.397 & 22.939 & 23.588 & 2.252 \\
1024 & 64   & 7.429 & 3.017 & 3.985 & 16.899 & 19.123 & 2.502 \\
1024 & 128  & 9.091 & 3.683 & 4.858 & 20.902 & 23.664 & 3.002 \\
1024 & 256  & 8.158 & 3.619 & 6.706 & 25.124 & 29.365 & 4.002 \\
1024 & 512  & 8.631 & 3.745 & 10.759 & 29.329 & 34.466 & 6.002 \\
1024 & 1024 & 7.758 & 3.697 & 16.398 & 47.715 & 62.911 & 10.002 \\
1024 & 2048 & 8.490 & 5.048 & 32.308 & 85.957 & 116.310 & 18.002 \\
\bottomrule
\end{longtable}
\end{landscape}

\begin{table}[ht]
\centering
\caption{Observed representation crossover widths.  A crossover is the first
tested \(p\) for which two-sided maintenance is no slower or no larger than
the named alternative.  ``None'' means no crossover in the tested range.}
\label{tab:e3-cross}
\begin{tabular}{@{}rrrrrr@{}}
\toprule
\(r\) & Drag & Drag+mat. & Fresh & Fresh+mat. & Memory \\
\midrule
256  & None & 512 & 1 & 1 & 128 \\
1024 & None & 512 & 1 & 1 & 512 \\
\bottomrule
\end{tabular}
\end{table}

\subsection{Results and implications}

The reviewer's thin-factor regime is real: when downstream computation can
retain factorized \(L^{-1}B\) and \(L^{-1}C\), drag-along remains faster at
every tested width, including \(p=2r\).  Conversely, fresh thin solves are
slower than the direct two-sided update even at \(p=1\).  If a dense operator
must be materialized, drag-along crosses behind the two-sided method at
\(p=512\) for both ranks.  The one-sided state reaches the dense-state memory
footprint at exactly \(p=r/2\): 128 for \(r=256\), and 512 for \(r=1024\).
All relative representation discrepancies are below
\(6.3\times10^{-7}\).

The experiment covers general \(M=BC^\top\) at fixed factor widths.  It does
not separately benchmark the symmetric specialization \(M=BB^\top\), nor a
single run in which \(p\) grows online.  The width sweep is sufficient for the
regime and memory crossover, but those unexecuted variants should not be
claimed as measured.

\paperchange{Add a regime table stating that one-sided propagation is
preferable when a fixed thin factor remains factorized throughout downstream
computation.  State that the proposed method is preferable when \(M\) is
dense, a dense \(\Aw\) is explicitly consumed, multiple changing factors must
be combined, or factor width grows enough to erase the memory advantage.
This concession directly resolves RVfB Q1 and makes the paper's scope more
credible.}

\section{E4: conditioning, near-duplicate kernels, and downdates}
\label{sec:e4}

\reviewer{URn7 Q1 asks whether numerical stability degrades for nearly
degenerate matrices and near the downdate boundary.}

\subsection{Controlled-spectrum positive updates}

We generate \(r=128\) SPD matrices with log-spaced spectra and initial target
condition numbers \(10^2,10^4,10^6,10^8\).  We set
\(M=G-\lambda_{\min}(G)I\), so \(\|\Aw\|_2\le1\), apply 200 positive rank-one
updates in FP32 or FP64, and compare no reset with fresh resets every 50 or
200 updates.  Exact state is accumulated in FP64.  Table~\ref{tab:e4-positive}
reports means over three seeds.  Because positive updates increase the small
eigenvalues, the condition number at step 200 can be substantially below the
initial target.

\begin{landscape}
\scriptsize
\begin{longtable}{@{}llrrrrrr@{}}
\caption{Controlled-spectrum stability.  ``Initial forward'' is the relative
\(\Aw\) error immediately after the initial factorization.  Values are means
over seeds; each ``max'' entry is the mean of the three per-seed maxima over
recorded checkpoints.  \(N=\infty\) denotes no reset.}
\label{tab:e4-positive}\\
\toprule
Initial \(\kappa\) & Prec. & \(N\) & Final \(\kappa\) & Initial forward &
Final forward & Max forward & Max operator backward \\
\midrule
\endfirsthead
\toprule
Initial \(\kappa\) & Prec. & \(N\) & Final \(\kappa\) & Initial forward &
Final forward & Max forward & Max operator backward \\
\midrule
\endhead
$10^{2}$ & FP32 & $\infty$ & 9.97e+01 & 1.11e-06 & 1.07e-05 & 1.07e-05 & 2.00e-05 \\
$10^{2}$ & FP32 & 50 & 9.97e+01 & 1.11e-06 & 1.11e-06 & 1.32e-06 & 1.29e-06 \\
$10^{2}$ & FP32 & 200 & 9.97e+01 & 1.11e-06 & 1.11e-06 & 5.89e-06 & 1.05e-05 \\
$10^{2}$ & FP64 & $\infty$ & 9.97e+01 & 0 & 4.93e-15 & 4.93e-15 & 2.98e-15 \\
$10^{2}$ & FP64 & 50 & 9.97e+01 & 0 & 0 & 3.18e-15 & 8.13e-16 \\
$10^{2}$ & FP64 & 200 & 9.97e+01 & 0 & 0 & 4.20e-15 & 2.26e-15 \\
\midrule
$10^{4}$ & FP32 & $\infty$ & 8.67e+03 & 4.21e-05 & 3.99e-05 & 4.21e-05 & 2.21e-05 \\
$10^{4}$ & FP32 & 50 & 8.67e+03 & 4.21e-05 & 3.80e-05 & 4.32e-05 & 1.19e-06 \\
$10^{4}$ & FP32 & 200 & 8.67e+03 & 4.21e-05 & 3.80e-05 & 4.21e-05 & 1.11e-05 \\
$10^{4}$ & FP64 & $\infty$ & 8.67e+03 & 0 & 1.43e-13 & 1.43e-13 & 2.64e-15 \\
$10^{4}$ & FP64 & 50 & 8.67e+03 & 0 & 0 & 1.28e-13 & 8.79e-16 \\
$10^{4}$ & FP64 & 200 & 8.67e+03 & 0 & 0 & 1.35e-13 & 1.92e-15 \\
\midrule
$10^{6}$ & FP32 & $\infty$ & 1.25e+05 & 2.48e-03 & 5.44e-04 & 2.48e-03 & 1.79e-05 \\
$10^{6}$ & FP32 & 50 & 1.25e+05 & 2.48e-03 & 5.39e-04 & 2.48e-03 & 1.30e-06 \\
$10^{6}$ & FP32 & 200 & 1.25e+05 & 2.48e-03 & 5.39e-04 & 2.48e-03 & 1.02e-05 \\
$10^{6}$ & FP64 & $\infty$ & 1.25e+05 & 0 & 1.59e-12 & 6.62e-12 & 2.99e-15 \\
$10^{6}$ & FP64 & 50 & 1.25e+05 & 0 & 0 & 6.62e-12 & 9.51e-16 \\
$10^{6}$ & FP64 & 200 & 1.25e+05 & 0 & 0 & 6.62e-12 & 2.20e-15 \\
\midrule
$10^{8}$ & FP32 & $\infty$ & 3.48e+05 & 1.56e-01 & 1.40e-03 & 1.56e-01 & 1.99e-05 \\
$10^{8}$ & FP32 & 50 & 3.48e+05 & 1.56e-01 & 1.40e-03 & 1.56e-01 & 1.20e-06 \\
$10^{8}$ & FP32 & 200 & 3.48e+05 & 1.56e-01 & 1.40e-03 & 1.56e-01 & 1.02e-05 \\
$10^{8}$ & FP64 & $\infty$ & 3.48e+05 & 0 & 4.25e-12 & 4.17e-10 & 2.94e-15 \\
$10^{8}$ & FP64 & 50 & 3.48e+05 & 0 & 0 & 4.17e-10 & 8.44e-16 \\
$10^{8}$ & FP64 & 200 & 3.48e+05 & 0 & 0 & 4.17e-10 & 2.02e-15 \\
\bottomrule
\end{longtable}
\end{landscape}

The positive-update pipeline remains backward stable: even at the most
extreme FP32 setting, operator backward error stays near \(2\times10^{-5}\)
without reset and near \(10^{-6}\) with 50-step resets.  Forward error,
however, is condition dependent.  At initial \(\kappa=10^8\), the FP32
whitened operator already has approximately 0.156 relative error before any
incremental update.  This is not instability caused by the orthogonal
transport; it reflects the conditioning of Cholesky factorization and the two
triangular solves that define whitening.  FP64 remains near machine precision.

\subsection{Near-duplicate RBF kernels}

We next construct rank-128 RBF kernel matrices in three dimensions, replace
one eighth of the inducing points by near duplicates at separations
\(\delta\in\{10^{-1},\ldots,10^{-6}\}\), and add \(10^{-6}\) jitter.  We test
100 positive updates with resets every 50 steps and three seeds.  The
\emph{cumulative} operator uses \(M=G-10^{-6}I\) and satisfies
\(\|\Aw\|_2\approx1\); the \emph{random} operator is normalized before
whitening but produces an unbounded \(\|\Aw\|_2\), intentionally testing the
regime outside the paper's clean bounded-operator lemma.

\begin{landscape}
\scriptsize
\begin{longtable}{@{}llrrrrrr@{}}
\caption{Near-duplicate RBF results, averaged over three seeds.  Final values
are recorded after the step-100 reset; ``max forward'' is the mean of the
three per-seed trajectory maxima.}
\label{tab:e4-kernel}\\
\toprule
Prec. & Operator & \(\delta\) & Final \(\kappa\) & \(\|\Aw\|_2\) &
Final forward & Max forward & Final operator backward \\
\midrule
\endfirsthead
\toprule
Prec. & Operator & \(\delta\) & Final \(\kappa\) & \(\|\Aw\|_2\) &
Final forward & Max forward & Final operator backward \\
\midrule
\endhead
FP32 & cumulative & 1.00e-01 & 4.51e+07 & 1.00e+00 & 8.48e-02 & 9.12e-02 & 1.75e-07 \\
FP32 & cumulative & 1.00e-02 & 4.73e+07 & 1.00e+00 & 1.03e-01 & 1.08e-01 & 1.80e-07 \\
FP32 & cumulative & 1.00e-03 & 4.73e+07 & 1.00e+00 & 1.02e-01 & 1.15e-01 & 1.84e-07 \\
FP32 & cumulative & 1.00e-04 & 4.73e+07 & 1.00e+00 & 1.02e-01 & 1.09e-01 & 1.82e-07 \\
FP32 & cumulative & 1.00e-05 & 4.73e+07 & 1.00e+00 & 9.46e-02 & 9.80e-02 & 1.75e-07 \\
FP32 & cumulative & 1.00e-06 & 4.73e+07 & 1.00e+00 & 8.16e-02 & 8.81e-02 & 1.64e-07 \\
\midrule
FP32 & random & 1.00e-01 & 4.51e+07 & 2.22e+05 & 6.58e-02 & 8.01e-02 & 3.20e-05 \\
FP32 & random & 1.00e-02 & 4.73e+07 & 3.11e+05 & 9.81e-02 & 1.13e-01 & 8.02e-04 \\
FP32 & random & 1.00e-03 & 4.73e+07 & 3.59e+05 & 1.10e-01 & 1.22e-01 & 7.23e-03 \\
FP32 & random & 1.00e-04 & 4.73e+07 & 3.61e+05 & 1.25e-01 & 1.28e-01 & 7.60e-04 \\
FP32 & random & 1.00e-05 & 4.73e+07 & 3.61e+05 & 1.19e-01 & 1.34e-01 & 1.63e-05 \\
FP32 & random & 1.00e-06 & 4.73e+07 & 3.61e+05 & 1.19e-01 & 1.29e-01 & 9.32e-06 \\
\midrule
FP64 & cumulative & 1.00e-01 & 4.51e+07 & 1.00e+00 & 0 & 2.61e-10 & 5.60e-16 \\
FP64 & cumulative & 1.00e-02 & 4.73e+07 & 1.00e+00 & 0 & 3.44e-10 & 5.54e-16 \\
FP64 & cumulative & 1.00e-03 & 4.73e+07 & 1.00e+00 & 0 & 3.44e-10 & 5.63e-16 \\
FP64 & cumulative & 1.00e-04 & 4.73e+07 & 1.00e+00 & 0 & 3.50e-10 & 5.64e-16 \\
FP64 & cumulative & 1.00e-05 & 4.73e+07 & 1.00e+00 & 0 & 3.54e-10 & 5.59e-16 \\
FP64 & cumulative & 1.00e-06 & 4.73e+07 & 1.00e+00 & 0 & 3.37e-10 & 5.42e-16 \\
\midrule
FP64 & random & 1.00e-01 & 4.51e+07 & 2.22e+05 & 0 & 2.68e-10 & 8.05e-14 \\
FP64 & random & 1.00e-02 & 4.73e+07 & 3.11e+05 & 0 & 3.49e-10 & 2.02e-12 \\
FP64 & random & 1.00e-03 & 4.73e+07 & 3.59e+05 & 0 & 4.21e-10 & 1.83e-11 \\
FP64 & random & 1.00e-04 & 4.73e+07 & 3.61e+05 & 0 & 4.41e-10 & 1.89e-12 \\
FP64 & random & 1.00e-05 & 4.73e+07 & 3.61e+05 & 0 & 4.36e-10 & 4.36e-14 \\
FP64 & random & 1.00e-06 & 4.73e+07 & 3.61e+05 & 0 & 4.12e-10 & 2.48e-14 \\
\bottomrule
\end{longtable}
\end{landscape}

The RBF experiment reinforces the controlled-spectrum result.  FP32 has
8--13\% forward error at \(\kappa\approx4.5\)--\(4.7\times10^7\), even
immediately after a reset, while the bounded cumulative operator retains
backward error near \(2\times10^{-7}\).  For arbitrary random \(M\),
\(\|\Aw\|_2\) grows to \(2.2\)--\(3.6\times10^5\); the bounded-operator
premise no longer holds, and FP32 operator backward error can reach
\(7.2\times10^{-3}\).  FP64 remains accurate in both cases.

\subsection{Near-boundary downdates}

For a downdate \(G^-=G-zz^\top\), define leverage
\(\ell=z^\top G^{-1}z<1\).  Appendix~C.4 predicts worst-direction
amplification \((1-\ell)^{-1}\).  We evaluate
\(\ell\in\{0.1,0.5,0.9,0.99,0.999\}\) at \(r\in\{64,256\}\), three seeds, and
compare the constructed worst direction with random FP32 roundoff.  The final
column gives the outer-product weight \(\beta\) in
\(G-\beta zz^\top\) that enforces a guard \(\beta\ell\le0.5\); it is not a
multiplier on \(z\) itself (the equivalent vector multiplier is
\(\sqrt{\beta}\)).

\begin{table}[ht]
\centering
\small
\caption{Downdate amplification.  Random-roundoff entries are mean \(\pm\)
sample standard deviation over three seeds.}
\label{tab:e4-down}
\begin{tabular}{@{}rrrrrr@{}}
\toprule
\(r\) & \(\ell\) & Bound & Worst observed & FP32 roundoff &
Guard weight \(\beta\) \\
\midrule
64 & 0.1 & 1.11 & 1.11 & \(1.00\pm0.0017\) & 1.0000 \\
64 & 0.5 & 2 & 2 & \(1.04\pm0.022\) & 1.0000 \\
64 & 0.9 & 10 & 10 & \(1.93\pm0.056\) & 0.5556 \\
64 & 0.99 & 100 & 100 & \(8.19\pm1.4\) & 0.5051 \\
64 & 0.999 & 1000 & 1000 & \(55.7\pm20\) & 0.5005 \\
\midrule
256 & 0.1 & 1.11 & 1.11 & \(1.00\pm0.00014\) & 1.0000 \\
256 & 0.5 & 2 & 2 & \(1.01\pm0.012\) & 1.0000 \\
256 & 0.9 & 10 & 10 & \(1.81\pm0.089\) & 0.5556 \\
256 & 0.99 & 100 & 100 & \(6.41\pm1.2\) & 0.5051 \\
256 & 0.999 & 1000 & 1000 & \(32.1\pm16\) & 0.5005 \\
\bottomrule
\end{tabular}
\end{table}

The observed worst-direction amplification matches the theoretical bound to
numerical precision at every rank and leverage.  Typical roundoff grows more
slowly than the worst case, but rises sharply as \(\ell\to1\).  Thus the
downdate warning is not merely theoretical.

\paperchange{State that the orthogonal transport is nonexpansive, but the full
pipeline inherits the conditioning of factorization and triangular solves.
Report backward error alongside forward error.  Recommend FP64, additional
jitter, or more frequent resets for extremely ill-conditioned whitening.
For downdates, expose leverage and either reject, damp to a specified threshold,
or refactorize as \(\ell\) approaches one.}

\section{E5: CUDA correctness, batching, and optimized fallback}
\label{sec:e5}

\reviewer{URn7 Q2 asks whether sequential Givens rotations limit GPU
parallelism and how performance scales with rank.}

\subsection{Method and algebraic equivalence}

On an H100 80\,GB GPU, we compare the Triton update with a fair PyTorch
baseline that forms \(G^+=G+zz^\top\), \(M^+=M+zw^\top\), computes a batched
Cholesky factorization, and performs two batched triangular solves.  CUDA
events measure median device time after warm-up.  The fused kernel is tested
at \(r\in\{64,128\}\) and independent-update batches
\(T\in\{1,4,16,64\}\).  A follow-up sweep replaces the original serial tiled
fallback at \(r\in\{256,512,1024\}\), \(T=1\), with a
baseline-compatible parallel path.  ``Baseline-compatible'' means that it
maintains exactly the same \(L\) and \(\Aw\) state as the original method; it
does not cache an inverse factor or change the mathematical output.  All
configurations were frozen before the ten-seed campaign; seed 0 supplied the
original measurements or configuration selection, and seeds 1--10 supply the
statistics below.

The main speedup comes from exposing parallelism already present in the
original algebra.  Let
\[
 u=L^{-1}z,\qquad y=L^{-1}w,\qquad
 S_i=1+\sum_{j=0}^{i}u_j^2,\quad S_{-1}=1.
\]
Because
\[
 G^+=L(I+uu^\top)L^\top,\qquad
 C=\operatorname{chol}(I+uu^\top),\qquad L^+=LC,
\]
the original sequence of Givens rotations has the closed-form coefficients
\[
 c_i=\sqrt{\frac{S_{i-1}}{S_i}},\qquad
 s_i=\frac{u_i}{\sqrt{S_i}}.
\]
The two vectors required by the rank-one operator correction are likewise
prefix expressions:
\[
 (C^{-1}u)_i=\frac{u_i}{\sqrt{S_{i-1}S_i}},\qquad
 (C^{-1}y)_i
 =c_i\left(y_i-\frac{u_i}{S_{i-1}}\sum_{j<i}u_jy_j\right).
\]
Thus the implementation obtains \(u\) and \(y\) in one two-right-hand-side
triangular solve and obtains every \(c_i,s_i,(C^{-1}u)_i,(C^{-1}y)_i\)
with parallel prefix sums.

The factor update is also a row-wise scan.  Writing
\(v_z=C^{-1}u\), for \(i\ge j\),
\[
 L^+_{ij}
 =\frac{L_{ij}}{c_j}
 +(v_z)_j\sum_{k=j+1}^{i}L_{ik}u_k.
\]
Each row therefore uses an independent reverse cumulative sum rather than an
\(r\)-step mutation chain.  Finally,
\[
 \Aw^+=C^{-1}\Aw C^{-\top}
       +(C^{-1}u)(C^{-1}y)^\top .
\]
Each left or right Givens sweep has the affine recurrence
\(p_i=c_i p_{i-1}-s_i x_i\), with output
\(x_i'=c_i x_i+s_i p_{i-1}\).  Associative composition of these affine maps
is an exact prefix scan in real arithmetic.  The CUDA path therefore
reorders, but does not approximate or omit, the original update; its only
difference is the usual FP32 reassociation error.  Reusable workspaces,
parallel row/column tiles, and fusion of the final outer product into the
right sweep remove the remaining allocator and launch overheads.

\begin{landscape}
{\scriptsize
\setlength{\tabcolsep}{1.5pt}
\begin{longtable}{@{}lrrrrrrrrrr@{}}
\caption{Ten-seed CUDA results on an H100 80\,GB.  Stochastic entries are
mean \(\pm\) sample SD across seed-level CUDA-event medians over 60 prepared
repeats.  Latency is microseconds and throughput is thousands of updates/s.
``BW'' is conservative minimum effective bandwidth rather than peak hardware
bandwidth.  ``Aux.'' is the exact, configuration-fixed sum of reusable
workspace and peak extra allocation, in KiB.}
\label{tab:e5}\\
\toprule
Impl. & \(r\) & \(T\) & Ours \(\mu\)s & Refact \(\mu\)s &
Ours k/s & Refact k/s & Speedup & BW GB/s & Aux. KiB &
\(\Aw\) rel.\ err. \\
\midrule
\endfirsthead
\toprule
Impl. & \(r\) & \(T\) & Ours \(\mu\)s & Refact \(\mu\)s &
Ours k/s & Refact k/s & Speedup & BW GB/s & Aux. KiB &
\(\Aw\) rel.\ err. \\
\midrule
\endhead
fused & 64 & 1  & \(105.1\pm0.5\) & \(125.0\pm1.0\) &
  \(9.52\pm0.05\) & \(8.00\pm0.06\) & \(1.190\pm0.010\) &
  \(0.643\pm0.003\) & 18.5 & \((2.42\pm0.17)\times10^{-7}\) \\
fused & 64 & 4  & \(105.7\pm0.6\) & \(222.1\pm3.2\) &
  \(37.9\pm0.2\) & \(18.0\pm0.3\) & \(2.10\pm0.03\) &
  \(2.558\pm0.014\) & 69 & \((1.86\pm0.04)\times10^{-7}\) \\
fused & 64 & 16 & \(105.8\pm0.4\) & \(199.8\pm1.4\) &
  \(151.2\pm0.5\) & \(80.1\pm0.6\) & \(1.888\pm0.011\) &
  \(10.22\pm0.04\) & 276 & \((2.187\pm0.020)\times10^{-7}\) \\
fused & 64 & 64 & \(106.3\pm0.5\) & \(201.4\pm1.0\) &
  \(602.2\pm2.7\) & \(317.8\pm1.6\) & \(1.895\pm0.015\) &
  \(40.7\pm0.2\) & 1,104 & \((2.191\pm0.006)\times10^{-7}\) \\
\midrule
fused & 128 & 1  & \(187.6\pm0.6\) & \(172.2\pm1.1\) &
  \(5.33\pm0.02\) & \(5.81\pm0.04\) & \(0.918\pm0.005\) &
  \(1.419\pm0.005\) & 66.5 & \((2.52\pm0.08)\times10^{-7}\) \\
fused & 128 & 4  & \(189.7\pm0.6\) & \(372.5\pm2.8\) &
  \(21.09\pm0.07\) & \(10.74\pm0.08\) & \(1.964\pm0.012\) &
  \(5.615\pm0.019\) & 266 & \((2.026\pm0.018)\times10^{-7}\) \\
fused & 128 & 16 & \(189.6\pm0.4\) & \(334.7\pm3.3\) &
  \(84.4\pm0.2\) & \(47.8\pm0.5\) & \(1.765\pm0.017\) &
  \(22.46\pm0.04\) & 1,064 & \((2.314\pm0.016)\times10^{-7}\) \\
fused & 128 & 64 & \(190.2\pm0.6\) & \(372.0\pm3.2\) &
  \(336.4\pm1.0\) & \(172.1\pm1.5\) & \(1.955\pm0.017\) &
  \(89.6\pm0.3\) & 4,256 & \((2.315\pm0.004)\times10^{-7}\) \\
\midrule
parallel & 256  & 1 & \(153.0\pm1.2\) & \(278.2\pm1.6\) &
  \(6.54\pm0.05\) & \(3.59\pm0.02\) & \(1.818\pm0.013\) &
  \(6.91\pm0.06\) & 521 & \((2.696\pm0.063)\times10^{-7}\) \\
parallel & 512  & 1 & \(265.4\pm1.9\) & \(505.0\pm1.5\) &
  \(3.77\pm0.03\) & \(1.980\pm0.006\) & \(1.903\pm0.014\) &
  \(15.87\pm0.11\) & 2,066 & \((2.612\pm0.025)\times10^{-7}\) \\
parallel & 1024 & 1 & \(496.9\pm3.5\) & \(1{,}023.7\pm2.6\) &
  \(2.013\pm0.014\) & \(0.977\pm0.002\) & \(2.060\pm0.013\) &
  \(33.83\pm0.24\) & 8,228 & \((3.040\pm0.015)\times10^{-7}\) \\
\bottomrule
\end{longtable}
}
\end{landscape}

All configurations are accurate: the ten-seed row-mean relative \(L\) errors
span \(0.90\)--\(1.40\times10^{-7}\), and the relative \(\Aw\) errors span
\(1.86\)--\(3.04\times10^{-7}\); the worst individual errors remain below
\(1.6\times10^{-7}\) and \(3.1\times10^{-7}\), respectively.  The fused path
demonstrates the value of independent-update batching.  At \(T=1\), \(r=128\)
is slightly slower than \refact{}; at \(T=4\)--64, fused mean speedups span
\(1.8\)--\(2.1\times\).
The baseline-compatible parallel fallback is also faster for every tested
large rank: \(1.8\times\), \(1.9\times\), and \(2.1\times\) at
\(r=256,512,1024\), respectively.  It replaces the old serial tiled fallback
without changing maintained state.  Separate seed-0-only phase profiling,
which is not part of the ten-seed timing summary, identifies the vendor
two-right-hand-side triangular solve as the remaining bottleneck; padding the
solve to wider right-hand sides did not improve it.

\subsection{Experimental inverse-factor state and refresh policy}
\label{sec:e5-inverse}

We separately tested a more aggressive representation that stores
\(L^{-1}\), rather than \(L\), between refreshes.  It uses the exact identity
\[
 (L^+)^{-1}=C^{-1}L^{-1},\qquad
 \bigl[(L^+)^{-1}\bigr]_{ij}
 =c_i\left((L^{-1})_{ij}
 -\frac{u_i}{S_{i-1}}\sum_{k<i}u_k(L^{-1})_{kj}\right),
\]
so future \(u\) and \(y\) are matrix--vector products and the inverse-factor
update is another column-wise prefix scan.  This preserves the update
algebra, but changes the stored factor representation and therefore requires
separate timing, drift, refresh, and memory evidence.

The first inverse-factor timing and stability jobs were invalidated because
\texttt{torch.linalg.inv} produced column-major-strided inputs while the
Triton kernel's contract is a row-major contiguous lower triangle.  Every
result below comes from corrected initialization that explicitly takes the
lower triangle and makes it contiguous.  Invalid jobs 393306 and 393301 remain
excluded.

\subsubsection{Held-out cached-state kernel timing}

The single-block configurations were selected using seed 0 and then frozen.
Job 393591 evaluated ten held-out problem seeds (1--10); each seed-level
latency is a CUDA-event median over 60 prepared repeats.
Table~\ref{tab:e5-inverse-heldout} reports the mean and sample standard
deviation across those ten seed-level medians.  Correctness cells use the
same ten-seed mean and sample standard deviation.

\begin{table}[ht]
\centering
\scriptsize
\caption{Corrected held-out single-block inverse-factor timing on one H100.
The incremental interval updates cached \(L^{-1}\) and \(\Aw\); initial inverse
construction, refresh, and maintenance of exact \(G,M\) are excluded.
\refact{} forms \(G^+,M^+\), factorizes, and rewhitens.  Thus this is an
isolated cached-state kernel comparison, not a self-contained application
measurement.  Stochastic entries are mean \(\pm\) sample SD over ten seeds.}
\label{tab:e5-inverse-heldout}
\begin{tabular}{@{}rcccrr@{}}
\toprule
\(r\) & Inverse \(\mu\)s & Refact \(\mu\)s & Speedup &
\(L^{-1}\) rel.\ err. & \(\Aw\) rel.\ err. \\
\midrule
256  & \(83.7\pm0.6\) & \(277.5\pm0.8\) &
  \(3.32\pm0.03\) & \((1.836\pm0.051)\times10^{-7}\) &
  \((2.696\pm0.063)\times10^{-7}\) \\
512  & \(83.2\pm0.6\) & \(506.4\pm2.1\) &
  \(6.09\pm0.06\) & \((1.805\pm0.015)\times10^{-7}\) &
  \((2.612\pm0.025)\times10^{-7}\) \\
1024 & \(83.9\pm0.7\) & \(1{,}025.0\pm4.8\) &
  \(12.22\pm0.12\) & \((2.051\pm0.011)\times10^{-7}\) &
  \((3.040\pm0.015)\times10^{-7}\) \\
\bottomrule
\end{tabular}
\end{table}

The held-out result establishes a large isolated-kernel advantage for a
consumer that already has the inverse-factor state: \(3.3\times\) at
\(r=256\), \(6.1\times\) at \(r=512\), and \(12.2\times\) at
\(r=1024\).  It does not by itself establish the amortized cost of constructing
or refreshing that state; the refresh and self-contained comparisons below
measure those costs separately.

\subsubsection{Repeated updates and condition dependence}

Job 393591 applies 500 updates without reset to ten well-conditioned
synthetic problems per row.  Table~\ref{tab:e5-inverse-stability} reports the
mean and sample standard deviation of the final error.  The first three rows
use the single-block
kernel; the last row independently checks the in-place recurrence used by the
large-rank path.

\begin{table}[ht]
\centering
\scriptsize
\caption{Corrected 500-update FP32 inverse-factor stability.  Stochastic
entries are mean \(\pm\) sample SD over ten seeds; the initial measured
condition numbers are near five.}
\label{tab:e5-inverse-stability}
{\setlength{\tabcolsep}{3pt}
\begin{tabular}{@{}lrrrrr@{}}
\toprule
Path & \(r\) & Initial \(\kappa(G)\) &
\(L^{-1}\) rel.\ err. & \(\Aw\) rel.\ err. & Inverse residual \\
\midrule
single-block & 256  & \(4.91\pm0.05\) &
  \((6.76\pm0.07)\times10^{-6}\) & \((1.193\pm0.011)\times10^{-5}\) &
  \((6.86\pm0.08)\times10^{-6}\) \\
single-block & 512  & \(4.96\pm0.04\) &
  \((1.018\pm0.009)\times10^{-5}\) & \((1.215\pm0.006)\times10^{-5}\) &
  \((1.068\pm0.011)\times10^{-5}\) \\
single-block & 1024 & \(4.97\pm0.03\) &
  \((7.66\pm0.04)\times10^{-6}\) & \((1.263\pm0.005)\times10^{-5}\) &
  \((7.77\pm0.05)\times10^{-6}\) \\
in-place     & 1024 & \(4.97\pm0.03\) &
  \((7.66\pm0.04)\times10^{-6}\) & \((1.450\pm0.005)\times10^{-5}\) &
  \((7.77\pm0.05)\times10^{-6}\) \\
\bottomrule
\end{tabular}
}
\end{table}

The well-conditioned errors grow gradually.  The largest \(\Aw\) error over
all ten seeds and all checkpoints is \(1.46\times10^{-5}\); this global
observed maximum is distinct from the mean \(\pm\) SD cells in the table.
This is not a condition-independent guarantee.  Job 393591 also fixes
\(r=256\), applies 200 updates, and varies the initial condition number.  As
Table~\ref{tab:e5-inverse-condition} shows, the operator error rises to about
\(4.06\pm0.16\%\) at \(\kappa(G)\approx10^6\), with a global observed maximum
of \(4.32\%\) across the ten seeds and checkpoints.  This is consistent with
the forward-conditioning mechanism isolated in E4.

\begin{table}[ht]
\centering
\scriptsize
\caption{Condition stress for the corrected single-block inverse-factor path
at \(r=256\) after 200 FP32 updates.  Stochastic entries are mean \(\pm\)
sample SD over ten seeds.}
\label{tab:e5-inverse-condition}
{\setlength{\tabcolsep}{3pt}
\begin{tabular}{@{}rrrrr@{}}
\toprule
Target \(\kappa(G)\) & Measured \(\kappa(G)\) &
\(L^{-1}\) rel.\ err. & \(\Aw\) rel.\ err. & Inverse residual \\
\midrule
\(10^2\) & \(99.999\pm0.001\) &
  \((3.02\pm0.07)\times10^{-6}\) & \((5.09\pm0.15)\times10^{-6}\) &
  \((3.07\pm0.04)\times10^{-6}\) \\
\(10^4\) & \(9{,}999.9\pm0.6\) &
  \((1.22\pm0.05)\times10^{-4}\) & \((1.82\pm0.11)\times10^{-4}\) &
  \((6.78\pm0.23)\times10^{-5}\) \\
\(10^6\) & \((1.003\pm0.004)\times10^6\) &
  \((2.32\pm0.07)\times10^{-2}\) & \((4.06\pm0.16)\times10^{-2}\) &
  \((8.96\pm0.25)\times10^{-3}\) \\
\bottomrule
\end{tabular}
}
\end{table}

\subsubsection{Refresh cost and drift}

Job 393591 measures 200-update sequences at refresh intervals
\(q\in\{10,50,100,200\}\), plus no refresh, for ten seeds at each rank.
Updates to exact \(G\) and \(M\) are inside every timed step.  At a refresh,
the timed region also includes Cholesky factorization, construction of
\(L^{-1}\) by a triangular solve, and complete rewhitening of \(M\).
\refact{} performs these operations every step.  Initial state construction is
outside both timed sequences.

Table~\ref{tab:e5-inverse-refresh} gives the mean and sample standard
deviation of CUDA-event time and paired speedup across seeds.  For each seed,
the error statistic is its maximum error immediately \emph{before} a refresh
(or at the final unrefreshed tail); the table then reports the mean and sample
standard deviation of those ten seed-level maxima, rather than the potentially
zero post-refresh final error.  The fresh-every-step baselines are
\(251.2\pm1.0\), \(478.8\pm2.1\), and \(988.7\pm8.3\)\,\(\mu\)s/update at
\(r=256,512,1024\), respectively.

\begin{table}[ht]
\centering
\caption{Periodic inverse-factor refresh over 200 updates.  Refresh work and
exact \(G,M\) maintenance are included in event timing.  \(q=\mathrm{none}\)
means no refresh.  Stochastic entries are mean \(\pm\) sample SD over ten
seeds; error entries aggregate each seed's maximum pre-refresh error.}
\label{tab:e5-inverse-refresh}
{\scriptsize
\setlength{\tabcolsep}{1.5pt}
\begin{tabular}{@{}rrrcrrr@{}}
\toprule
\(r\) & \(q\) & Refreshes & Event \(\mu\)s/update & Speedup &
Seed-max \(L^{-1}\) err. & Seed-max \(\Aw\) err. \\
\midrule
256 & none & 0  & \(100.4\pm0.8\) & \(2.50\pm0.02\) &
  \((2.76\pm0.06)\times10^{-6}\) & \((4.84\pm0.09)\times10^{-6}\) \\
256 & 10   & 20 & \(120.8\pm0.8\) & \(2.08\pm0.02\) &
  \((2.76\pm0.03)\times10^{-7}\) & \((4.45\pm0.06)\times10^{-7}\) \\
256 & 50   & 4  & \(103.1\pm1.1\) & \(2.44\pm0.03\) &
  \((7.77\pm0.14)\times10^{-7}\) & \((1.344\pm0.029)\times10^{-6}\) \\
256 & 100  & 2  & \(100.3\pm0.8\) & \(2.51\pm0.02\) &
  \((1.440\pm0.026)\times10^{-6}\) & \((2.51\pm0.05)\times10^{-6}\) \\
256 & 200  & 1  & \(98.7\pm1.0\) & \(2.54\pm0.03\) &
  \((2.76\pm0.06)\times10^{-6}\) & \((4.84\pm0.09)\times10^{-6}\) \\
\midrule
512 & none & 0  & \(101.3\pm1.6\) & \(4.73\pm0.07\) &
  \((4.11\pm0.05)\times10^{-6}\) & \((4.93\pm0.05)\times10^{-6}\) \\
512 & 10   & 20 & \(135.3\pm0.9\) & \(3.54\pm0.03\) &
  \((3.22\pm0.04)\times10^{-7}\) & \((4.35\pm0.05)\times10^{-7}\) \\
512 & 50   & 4  & \(107.7\pm1.2\) & \(4.45\pm0.06\) &
  \((1.097\pm0.012)\times10^{-6}\) & \((1.355\pm0.021)\times10^{-6}\) \\
512 & 100  & 2  & \(102.3\pm0.7\) & \(4.68\pm0.04\) &
  \((2.098\pm0.027)\times10^{-6}\) & \((2.545\pm0.030)\times10^{-6}\) \\
512 & 200  & 1  & \(102.7\pm4.7\) & \(4.67\pm0.19\) &
  \((4.11\pm0.05)\times10^{-6}\) & \((4.93\pm0.05)\times10^{-6}\) \\
\midrule
1024 & none & 0  & \(100.6\pm1.1\) & \(9.83\pm0.13\) &
  \((3.106\pm0.015)\times10^{-6}\) & \((5.120\pm0.028)\times10^{-6}\) \\
1024 & 10   & 20 & \(198.3\pm1.3\) & \(4.99\pm0.05\) &
  \((2.986\pm0.020)\times10^{-7}\) & \((4.659\pm0.051)\times10^{-7}\) \\
1024 & 50   & 4  & \(116.0\pm0.5\) & \(8.52\pm0.08\) &
  \((8.566\pm0.070)\times10^{-7}\) & \((1.404\pm0.013)\times10^{-6}\) \\
1024 & 100  & 2  & \(117.0\pm9.0\) & \(8.49\pm0.64\) &
  \((1.600\pm0.013)\times10^{-6}\) & \((2.632\pm0.022)\times10^{-6}\) \\
1024 & 200  & 1  & \(105.2\pm0.6\) & \(9.40\pm0.09\) &
  \((3.106\pm0.015)\times10^{-6}\) & \((5.120\pm0.028)\times10^{-6}\) \\
\bottomrule
\end{tabular}
}
\end{table}

Every tested refresh policy is faster than refactorizing every step in all ten
seeds; the minimum paired speedup over all rows and seeds is \(2.04\times\).
Refreshing every 10 steps holds the global observed pre-refresh \(\Aw\) error below
\(4.76\times10^{-7}\) across all ten seeds and ranks, but pays the highest
refresh cost.  An interval of 50 steps holds the corresponding global
observed maximum below \(1.43\times10^{-6}\), while retaining mean speedups of
\(2.44\pm0.03\times\), \(4.45\pm0.06\times\), and
\(8.52\pm0.08\times\) at the three ranks.  These global maxima are safety
summaries distinct from the mean \(\pm\) SD seed-max cells in the table.
These are controlled synthetic sequences, so \(q=50\) is an
illustrative cost--drift point, not a universal policy; conditioning and the
application's error tolerance should determine the reset interval.  This
numerical refresh is distinct from the chart-refresh threshold \(\tref\) in
E1.

The \(q=50\) timing is stable across seeds at every rank (event-time SD at
most \(1.2\,\mu\)s).  By contrast, the \(r=1024,q=100\) row is bimodal
(\(117.0\pm9.0\,\mu\)s and \(8.49\pm0.64\times\)), and the
\(r=512,q=200\) row contains one process-level timing outlier.  We therefore
do not interpret \(q=100\) as a stable timing point; the \(q=50\)
recommendation does not depend on either anomalous row.

\subsubsection{Synthetic self-contained large-rank comparison}

The full-axis in-place path removes the single-block rank limit while retaining
exact \(G,M\) for future refresh.  The comparison maintains four dense FP32
state matrices on both sides: \(G,M,L^{-1},\Aw\) for the inverse-factor path
and \(G,M,L,\Aw\) for \refact{}.  Every timed event includes the \(G,M\)
rank-one outer updates.  Each repeat prepares the same well-scaled state
outside timing and measures one recurrent-compatible update, so this is not a
long-horizon stability experiment.  The ten seeds generate independent update
vectors.  To counterbalance order effects, odd seeds measure the inverse path
first and even seeds measure refactorization first at every rank.

Tile and warp choices were selected at \(r=8192\) and \(r=32768\) in job
393466; none of those tuning latencies enters
Table~\ref{tab:e5-large-scale}.  The final ten-seed measurement ladder is
job 393591.
The implementation is compared numerically with a fresh factorization through
\(r=32{,}768\).  At larger ranks, memory-safe validation checks every output
element for finiteness but does not compute a fresh numerical error.

\begin{table}[ht]
\centering
\scriptsize
\caption{Final self-contained large-rank ladder on one H100 80\,GB.
Stochastic entries are mean \(\pm\) sample SD across ten seed-level
CUDA-event medians; within-seed repeat counts are 10 through \(r=8192\), 5
through \(r=32768\), and 3 thereafter.  Both methods maintain four dense
matrices; ``State'' is exact minimum FP32 storage.  Fresh-check entries report
mean \(\pm\) sample SD for \(L^{-1}/\Aw\) relative errors.}
\label{tab:e5-large-scale}
{\setlength{\tabcolsep}{2.5pt}
\begin{tabular}{@{}rrrrrl@{}}
\toprule
\(r\) & Inverse ms & Refact ms & Speedup & State GiB & Validation \\
\midrule
2,048  & \(1.432\pm0.007\) & \(2.550\pm0.012\) &
  \(1.781\pm0.008\) & 0.0625 &
  fresh: \((3.94\pm0.08)\times10^{-8}/(9.72\pm0.19)\times10^{-8}\) \\
4,096  & \(3.550\pm0.006\) & \(7.694\pm0.036\) &
  \(2.167\pm0.008\) & 0.25 &
  fresh: \((4.93\pm0.06)\times10^{-8}/(1.215\pm0.040)\times10^{-7}\) \\
8,192  & \(7.933\pm0.011\) & \(39.7\pm0.1\) &
  \(5.00\pm0.01\) & 1 &
  fresh: \((3.83\pm0.04)\times10^{-8}/(9.95\pm0.11)\times10^{-8}\) \\
16,384 & \(19.85\pm0.01\) & \(246.1\pm0.8\) &
  \(12.39\pm0.04\) & 4 &
  fresh: \((4.61\pm0.03)\times10^{-8}/(1.263\pm0.057)\times10^{-7}\) \\
32,768 & \(58.0\pm0.1\) & \(1{,}755.1\pm2.6\) &
  \(30.29\pm0.05\) & 16 &
  fresh: \((3.416\pm0.015)\times10^{-8}/(1.049\pm0.014)\times10^{-7}\) \\
49,152 & \(113.0\pm0.1\) & \(5{,}717.1\pm4.2\) &
  \(50.61\pm0.06\) & 36 & 10/10 full finite \\
65,536 & \(183.6\pm0.2\) & \(13{,}093.8\pm9.3\) &
  \(71.3\pm0.1\) & 64 & 10/10 full finite \\
\bottomrule
\end{tabular}
}
\end{table}

Every paired seed is already favorable at \(r=2048\), and the mean speedup
grows with rank as the quadratic inverse-factor update is compared with cubic
factorization and dense triangular solves.  At \(r=65{,}536\), the final event
times are \(183.6\pm0.2\)\,ms versus \(13.094\pm0.009\)\,s, a
\(71.3\pm0.1\times\) paired speedup.  This
largest row is a full finite-output check, not a fresh numerical-correctness
check, and the 64\,GiB four-matrix state leaves limited headroom on the
80\,GB device.  We therefore make no claim beyond the measured range.

\subsubsection{Real-data large-rank scaling on The Well}
\label{sec:e5-blackhole}

To answer URn7 Q2's request to analyze scaling at large rank, we froze a
real-data extension before feature generation or benchmark timing.  This is a
streaming Koopman/operator-regression workload, not a predictive-accuracy or
natural-gradient experiment.  We use the official Polymathic AI
\href{https://huggingface.co/datasets/polymathic-ai/post_neutron_star_merger}%
{\texttt{post\_neutron\_star\_merger}} dataset from The Well, training
scenario~4 (\texttt{torus\_MBH\_10}), a disk from a black-hole--neutron-star
merger with a 10-solar-mass black hole.  The pinned Hugging Face revision is
\texttt{1fd582f3e29749ac4d0c14bc816bddb2394258cc}.  The exact
14,109,638,656-byte file contains one 181-snapshot trajectory on a
\(192\times128\times66\) spherical grid.

The feature map uses twelve dynamic channels: density, internal energy,
pressure, temperature, electron fraction, entropy, and the three spatial
components of both magnetic field and velocity.  We exclude the
time-independent contravariant spacetime metric.  Per-channel scalar means
and population standard deviations are accumulated in FP64 over every
spatial point in snapshots 0--63 and then fixed.  Each standardized channel
receives an orthonormal separable 3-D DCT-II.  Spatial frequencies are ordered
by
\[
 (k_{\log r}/192)^2+(k_\theta/128)^2+(k_\phi/66)^2
\]
with ascending lexicographic indices as the exact tie-break; coefficients are
interleaved frequency-major and channel-minor.  This yields one C-contiguous
FP32 trace of shape \(181\times65{,}536\), and every tested representation is
an exact prefix of that trace.  For rank \(r\), the unscaled prefix is divided
by
\[
 s_r^2=\frac1{64}\sum_{t=0}^{63}\|\phi_r(x_t)\|_2^2 ,
\]
so feature energy is comparable across ranks without using an evaluation
transition.

Writing \(z_t=\phi_r(x_t)/s_r\) and
\(w_t=\phi_r(x_{t+1})/s_r\), the common state is
\[
 G_0=I+\sum_{t=0}^{63}z_tz_t^\top,\qquad
 M_0=\sum_{t=0}^{63}z_tw_t^\top,\qquad
 {\Aw}_0=L_0^{-1}M_0L_0^{-\top};
\]
the ridge coefficient is one.  Transition 64 is used only for compilation
and warm-up.  The ten predeclared evaluation transitions are
\(80,91,102,113,124,135,146,157,168,179\).  Each is applied separately to
the same state after transition 63, isolating one recurrent update rather
than a long-horizon drift trajectory.  Even-numbered cases execute the
inverse path first and odd-numbered cases reverse the method order.

Both methods maintain four dense FP32 matrices:
\(G,M,L^{-1},\Aw\) for the experimental inverse-factor path and
\(G,M,L,\Aw\) for \refact{}.  The inverse kernel consumes a row-major
contiguous, explicitly lower-triangular \(L^{-1}\).  It packs \(z,w\) into
one two-column right-hand side and obtains
\(u=L^{-1}z,\ y=L^{-1}w\) with one dense matrix multiplication.
Elementwise \(u_i^2\) and \(u_i y_i\) terms are prefix-summed to form the
exact Givens parameters and correction vectors.  One Triton program owns
each independent column tile, updates \(L^{-1}\) in place, and applies the
left sweep to \(\Aw\); a second owns independent row tiles and fuses the
right sweep with the rank-one \(v_zv_w^\top\) correction.  Each recurrence
remains serial within its row or column, while independent chains execute
concurrently.  The kernel therefore performs \(O(r^2)\) work without an
\(r\times r\) scratch matrix; its reusable workspace is \(9r\) FP32 values
(five vectors and two \(r\times2\) buffers).  Frozen row/column tiles are
\(8,8,8,16,32,32,32\) at the seven ascending ranks; setup uses four warps
and each sweep one.  \refact{} uses CUDA \texttt{cholesky\_ex} followed by
the two triangular solves that form the newly whitened operator.  TF32 is
disabled.

A CUDA event includes the common dense updates
\(G\leftarrow G+zz^\top\) and \(M\leftarrow M+zw^\top\), followed by the
selected method core.  Feature I/O and DCT work, host-to-device transfer,
initial-state construction, state restoration, compilation, warm-up,
pre-timing synchronization, and correctness checks are excluded.  Each
transition contributes its within-case event median: ten repeats through
\(r=8192\), five through \(r=32768\), and three thereafter.  Paired speedup
is formed within transition, then summarized by its mean, sample SD, and
two-sided 95\% Student-\(t\) interval across the ten transitions.  Timing
repeats are retained in the raw JSON but are not pseudoreplicated as
independent observations, and no timing-based filtering is performed.

\begin{table}[ht]
\centering
\scriptsize
\caption{Locked real-data rank scaling on one H100 per row.  Event times are
mean \(\pm\) sample SD in milliseconds; speedup is the mean \(\pm\) sample SD
of the ten paired transition ratios.  The interval is the two-sided 95\%
\(t_9\) interval for paired speedup.  ``State'' is the exact minimum storage
for four dense FP32 matrices.}
\label{tab:e5-blackhole-scale}
{\setlength{\tabcolsep}{2.6pt}
\begin{tabular}{@{}rrrrrr@{}}
\toprule
\(r\) & Inverse ms & Refact ms & Paired speedup &
95\% interval & State GiB \\
\midrule
2,048  & \(1.494\pm0.007\) & \(2.5578\pm0.0029\) &
  \(1.7119\pm0.0085\) & \([1.7058,1.7179]\) & 0.0625 \\
4,096  & \(3.576\pm0.021\) & \(7.6513\pm0.0025\) &
  \(2.1395\pm0.0120\) & \([2.1310,2.1481]\) & 0.25 \\
8,192  & \(7.948\pm0.014\) & \(39.7979\pm0.0141\) &
  \(5.0076\pm0.0104\) & \([5.0002,5.0151]\) & 1 \\
16,384 & \(19.792\pm0.012\) & \(247.008\pm0.016\) &
  \(12.4802\pm0.0074\) & \([12.4749,12.4855]\) & 4 \\
32,768 & \(58.093\pm0.103\) & \(1{,}755.897\pm0.337\) &
  \(30.2255\pm0.0577\) & \([30.1842,30.2668]\) & 16 \\
49,152 & \(113.944\pm0.182\) & \(5{,}696.014\pm8.706\) &
  \(49.9896\pm0.1153\) & \([49.9071,50.0721]\) & 36 \\
65,536 & \(183.720\pm0.188\) & \(13{,}182.753\pm33.306\) &
  \(71.7546\pm0.2247\) & \([71.5939,71.9153]\) & 64 \\
\bottomrule
\end{tabular}
}
\end{table}

\begin{figure}[ht]
\centering
\includegraphics[width=\linewidth]{%
  blackhole/blackhole-s4-v1-locked-rerun2/rank_scaling.pdf}
\caption{The Well isolated-transition ladder.  Error bars are 95\%
Student-\(t\) intervals across the ten predeclared transition medians, not
repeat-level intervals.  Both methods include the common \(G,M\) outer
updates.}
\label{fig:e5-blackhole-scale}
\end{figure}

Correctness is evaluated after timing.  Every fresh factorization must return
Cholesky \texttt{info=0}, and every element of the state and output must be
finite.  At every rank, sixteen normalized Rademacher probes with fixed seed
29421 compare
\[
 L_{\rm inc}^{-1}P\quad\text{with}\quad
 \operatorname{solve}(L_{\rm fresh},P),\qquad
 {\Aw}_{\rm inc}P\quad\text{with}\quad{\Aw}_{\rm fresh}P,
\]
and independently check the factor and whitened-operator identities.  Full
dense relative errors against fresh \(L^{-1}\) and \(\Aw\) are additionally
retained through \(r=32{,}768\).
Ranks 49,152 and 65,536 retain the fresh action-probe comparison and a
full-element finiteness scan, but not a full dense difference.  A row is
eligible for a speedup claim only if every predeclared relative error is at
most \(10^{-3}\).  Because each updated metric is \(I+ZZ^\top\) with only 65
stored vectors and \(r>65\), \(\lambda_{\min}(G)=1\) exactly; we compute
\(\kappa_2(G)=1+\lambda_{\max}(Z^\top Z)\) from the \(65\times65\) FP64
Gram matrix.

\begin{table}[ht]
\centering
\scriptsize
\caption{Real-data correctness over ten transitions per rank.  Probe columns
are the maximum cross-method relative error; the dense column is the maximum
direct incremental-versus-fresh \(L^{-1}/\Aw\) relative-Frobenius error.
Every predeclared residual, including diagnostics not shown, passed the
\(10^{-3}\) gate.}
\label{tab:e5-blackhole-valid}
{\setlength{\tabcolsep}{2.2pt}
\begin{tabular}{@{}rcrrrl@{}}
\toprule
\(r\) & \(\kappa_2(G)\) range &
\(\max L^{-1}P\) & \(\max \Aw P\) &
Dense \(L^{-1}/\Aw\) & Scope \\
\midrule
2,048  & 57.69--57.99 & \(2.69\times10^{-7}\) & \(4.18\times10^{-5}\) &
  \(2.33\times10^{-7}/1.53\times10^{-5}\) & dense+probe+finite \\
4,096  & 56.45--56.75 & \(2.10\times10^{-7}\) & \(4.59\times10^{-5}\) &
  \(1.76\times10^{-7}/1.65\times10^{-5}\) & dense+probe+finite \\
8,192  & 55.12--55.41 & \(1.62\times10^{-7}\) & \(3.44\times10^{-5}\) &
  \(1.33\times10^{-7}/1.81\times10^{-5}\) & dense+probe+finite \\
16,384 & 53.73--54.02 & \(1.34\times10^{-7}\) & \(2.73\times10^{-5}\) &
  \(1.06\times10^{-7}/1.80\times10^{-5}\) & dense+probe+finite \\
32,768 & 52.36--52.64 & \(1.11\times10^{-7}\) & \(2.62\times10^{-5}\) &
  \(8.71\times10^{-8}/2.46\times10^{-5}\) & dense+probe+finite \\
49,152 & 51.57--51.84 & \(1.13\times10^{-7}\) & \(2.71\times10^{-5}\) &
  -- & probe+finite \\
65,536 & 51.01--51.28 & \(1.04\times10^{-7}\) & \(2.55\times10^{-5}\) &
  -- & probe+finite \\
\bottomrule
\end{tabular}
}
\end{table}

Table~\ref{tab:e5-blackhole-scale} and
Figure~\ref{fig:e5-blackhole-scale} show that all seven rows and all 70 paired
transitions favor the inverse-factor path; Table~\ref{tab:e5-blackhole-valid}
confirms that every row passes the frozen correctness gate.  Mean paired
speedup grows monotonically
from \(1.7119\pm0.0085\times\) at \(r=2048\) to
\(71.7546\pm0.2247\times\) at \(r=65{,}536\); at the latter rank, mean event
times are \(183.720\pm0.188\)\,ms and
\(13.182753\pm0.033306\)\,s, respectively.  The largest of any predeclared
relative-error diagnostics is \(1.84\times10^{-4}\), below the \(10^{-3}\)
gate.  The increasing advantage is consistent with the common \(O(r^2)\)
outer updates plus an \(O(r^2)\) inverse-factor core versus dense
\(O(r^3)\) factorization and solves.  This independently extends the
synthetic ladder to structured update vectors from a physical trajectory; it
does not establish predictive accuracy or remove the inverse-state refresh
requirement.

The statistical interval is descriptive across ten temporally related
transitions from one trajectory, not across independent simulations.  Each
rank used a distinct H100, so paired comparisons are within one GPU whereas
the cross-rank curve also contains device-to-device variation.  This is an
isolated restored-state update in FP32, not end-to-end HDF5/DCT throughput or
a large-rank repeated-update stability test.  No specialized Woodbury
baseline is included, and the four matrices require 64\,GiB at the largest
rank.  The black-hole dataset was frozen as the primary real-data result
regardless of favorability; it was not selected after comparison with HIGGS.
We therefore do not extrapolate beyond \(65{,}536\).

\paperchange{Use the fused batching results and the baseline-compatible
parallel rows of Table~\ref{tab:e5} to answer URn7 Q2, and explain the exact
prefix-scan reformulation.  Present Tables~\ref{tab:e5-inverse-heldout}--%
\ref{tab:e5-blackhole-valid} separately as evidence for an experimental
inverse-factor state: distinguish isolated cached-state timing from
self-contained timing, report the condition-dependent drift and measured
refresh trade-off, add the frozen real-trajectory ladder, and limit
large-rank claims to each table's stated validation scope.}

\section{E6: practical \texorpdfstring{rank-\(k\)}{rank-k} cost}
\label{sec:e6}

\reviewer{URn7 Q4 asks about constants when extending rank-one updates to
rank-\(k\) changes.}

\subsection{Motivation and method}

For \(Z,W\in\R^{r\times k}\), we compare \(k\) sequential calls to the FP64
single-thread rank-one CPU kernel with one fresh factorization of
\[
  G^+=G+ZZ^\top,\qquad M^+=M+ZW^\top.
\]
We test \(r\in\{256,1024\}\), \(k\in\{1,2,4,8,16,32\}\), use median latency
over 20 repetitions (10 for \(k>8\)), and verify against fresh FP64
factorization.  No native blocked rank-\(k\) implementation is benchmarked.

\begin{table}[ht]
\centering
\small
\caption{Practical rank-\(k\) CPU results.  Latencies are milliseconds.
Speedup is \refact{} time divided by sequential rank-one time.}
\label{tab:e6}
\begin{tabular}{@{}rrrrrr@{}}
\toprule
\(r\) & \(k\) & Sequential & Refact & Speedup & \(\Aw\) rel.\ error \\
\midrule
256 & 1  & 0.418 & 1.150 & 2.751 & 5.12e-16 \\
256 & 2  & 0.813 & 1.134 & 1.394 & 5.40e-16 \\
256 & 4  & 1.635 & 1.139 & 0.697 & 5.62e-16 \\
256 & 8  & 3.212 & 1.133 & 0.353 & 6.03e-16 \\
256 & 16 & 6.343 & 1.196 & 0.189 & 6.97e-16 \\
256 & 32 & 12.642 & 1.210 & 0.096 & 8.43e-16 \\
\midrule
1024 & 1  & 7.776 & 70.548 & 9.072 & 6.20e-16 \\
1024 & 2  & 19.908 & 84.001 & 4.219 & 6.33e-16 \\
1024 & 4  & 44.213 & 92.219 & 2.086 & 6.56e-16 \\
1024 & 8  & 90.744 & 74.964 & 0.826 & 6.98e-16 \\
1024 & 16 & 159.037 & 77.950 & 0.490 & 7.74e-16 \\
1024 & 32 & 303.179 & 75.154 & 0.248 & 9.14e-16 \\
\bottomrule
\end{tabular}
\end{table}

\subsection{Results and implications}

Sequential rank-one updates are exact to FP64 precision and valuable for small
\(k\).  At \(r=256\), the measured crossover lies between \(k=2\) and \(k=4\);
at \(r=1024\), it lies between \(k=4\) and \(k=8\).  This is much earlier than
the FLOP-only predictions \(k\approx33.2\) and \(k\approx132.7\),
respectively.  The discrepancy reflects optimized cubic BLAS/LAPACK kernels,
cache behavior, and repeated rank-one call overhead.  It does not contradict
the asymptotic cost, but it invalidates a universal practical crossover
derived solely from flop counts.  A blocked compact-WY/Householder or blocked
Givens implementation may move the crossover, but was not run.

\paperchange{State the complete rank-\(k\) complexity
\(O(r^2k+rk^2+k^3)\), reducing to \(O(r^2k)\) only for \(k\ll r\).  Report
the strong \(k=1\) speedups and the measured repeated-update crossover.
Present the \(k\approx0.13r\) calculation only as an arithmetic model, not a
wall-clock prediction, and identify a native blocked rank-\(k\) kernel as
future work.}

\section{Integrated manuscript and rebuttal recommendations}
\label{sec:recommendations}

\subsection{Reviewer RVfB}

\begin{itemize}
  \item \textbf{Q1, one-sided alternatives.} Agree explicitly that a fixed,
  thin factor retained in factorized form is a regime where one-sided
  propagation is preferable.  Cite the absence of a drag-only crossover in
  Table~\ref{tab:e3-cross}.  Explain that direct two-sided maintenance is
  advantageous when an explicit dense operator is required: it beats fresh
  solves at the first tested width and beats drag-plus-materialization by
  \(p=512\).
  \item \textbf{Q2, predictive accuracy.} Withdraw the claim that the exact
  update itself improves test MSE.  Report the leakage-controlled E2
  time-ceiling rerun in Table~\ref{tab:e2-corrected}: at matched
  \(r=1024\), the paired \ours{}--\refact{} difference is
  \(0.005079\pm0.005808\), with no accuracy-win evidence.  Retain the E1
  matched-rank rows as independent supporting evidence, and do not describe
  the lower-rank E2 rows as equal-compute because their ten-epoch cap leaves
  most of the ceiling unused.
\end{itemize}

\subsection{Reviewer fdru}

\begin{itemize}
  \item \textbf{Q1, refresh threshold.} Report the smooth refresh-frequency
  curve from E1.  Use the refresh counts and MSE intervals to suggest
  \(\tref\approx0.05\)--0.1 as a starting range, while requiring
  system-specific timing calibration rather than treating the E1 seconds as
  portable.  Clarify that this threshold is application-specific and distinct
  from a periodic numerical reset.
  \item \textbf{Q2, dissemination.} Give an anonymized release plan with C++,
  Triton, and MLX build instructions, correctness tests, supported
  software/hardware, and a public-release date.  Treat AMD and Intel kernels as
  future work, consistent with the reviewer's statement that they are not
  required.
  \item \textbf{Notation.} Introduce \(z\) as the metric-update vector and
  \(w\) as the operator-update vector before Algorithm~1.
\end{itemize}

\subsection{Reviewer URn7}

\begin{itemize}
  \item \textbf{Q1, stability.} Distinguish forward conditioning from
  backward stability.  The orthogonal transport is stable, FP64 is robust,
  and resets control accumulated backward error; FP32 whitening can still
  have large forward error when \(\kappa(G)\) is extreme.  Expose and guard
  downdate leverage.
  \item \textbf{Q2, GPU scaling.} Report fused FP32 correctness and batching
  speedups together with the tuned, baseline-compatible prefix-scan fallback.
  Separately report the corrected held-out inverse-factor timing and synthetic
  condition-dependent 500-update drift and refresh-cost sweep.  Add the
  provenance-locked The Well isolated-transition ladder through
  \(r=65{,}536\), while keeping this representation explicitly experimental,
  distinguishing full-dense checks from fresh action probes, and stating its
  refresh and four-matrix memory requirements.
  \item \textbf{Q3, inverse-free methods.} Add a textual regime comparison:
  inverse-free methods are attractive for batch-oriented, matrix-multiply
  workloads; the proposed method is for systems that already maintain a
  sequential Cholesky state and consume whitened quantities.
  \item \textbf{Q4, rank-\(k\).} Add the full complexity and the measured
  small-\(k\) crossover from E6.
  \item \textbf{Q5, whitening.} Explain that a non-whitened model pays a
  one-time coordinate transformation; maintenance is worthwhile only when
  downstream computation repeatedly consumes the whitened operator.
\end{itemize}

\subsection{Area Chair}

The integrated response should emphasize six concrete changes:
(i) a one-sided/two-sided regime table;
(ii) corrected fixed-rank MSE wording, supported by the matched-rank
Table~\ref{tab:e1} rows and the leakage-controlled matched-rank result in
Table~\ref{tab:e2-corrected}, while explicitly withdrawing the old
checkpoint-truncated comparison;
(iii) the E1 threshold sweep;
(iv) forward and backward stability evidence plus a downdate guard;
(v) supported fused CUDA batching and baseline-compatible parallel-fallback
results, with the corrected inverse-factor timing, refresh, and large-rank
evidence---including the frozen real-trajectory ladder---separated as an
experimental stored-state representation;
and
(vi) the complete rank-\(k\) complexity with measured constants.
Independently, the abstract and introduction should state explicitly that
\(G,M,\Aw\in\R^{r\times r}\), \(G\) is SPD, and \(M\) may be arbitrary.

\section{Conclusion}

The reviewer-directed experiments support the paper's central message while
substantially improving its precision.  They confirm exactness and numerical
agreement, identify a practical threshold range, validate the one-sided
regime raised by RVfB, demonstrate backward stability and exact downdate
amplification, and show that fused CUDA batching is effective.  They also
show that the exact prefix-scan fallback is faster than CUDA refactorization
through rank 1024.  The corrected inverse-factor follow-up adds a
condition-aware refresh policy and a synthetic self-contained measured ladder
from \(r=2048\) to \(r=65{,}536\).  The independently locked real-data
extension uses structured transitions from The Well and reaches
\(71.7546\pm0.2247\times\) at \(r=65{,}536\), with a 95\% interval of
\([71.5939,71.9153]\), while all seven ranks pass their frozen correctness
gate.  It also identifies the limits of that result: the representation
changes stored state, drift grows sharply under extreme conditioning, the
real-data intervals describe ten related transitions from one trajectory,
and full-dense numerical comparison stops at \(r=32{,}768\) (fresh action
probes continue to \(65{,}536\)).  More broadly,
the corrected E2 matched-rank comparison gives a paired
\ours{}--\refact{} MSE difference of \(0.005079\pm0.005808\) and provides no
evidence that the update itself improves MSE,
extreme FP32 whitening is not condition-number independent, and practical
rank-\(k\) crossover is not justified by asymptotic flop counts alone.
Incorporating these distinctions would make the revision more credible,
easier to apply, and more directly responsive to all reviewer questions.

\appendix

\section{Result provenance and regeneration}

The tables in this report are derived from the planned-task outputs under
\path{rebuttal_results/}.  The canonical machine-readable artifact is the
curated set of individual JSON files, which retain seeds, repeat counts and,
where recorded, raw timing samples, environment metadata, full training
histories, and checkpoint diagnostics.  The task definitions are in
\path{rebuttal_experiments/task_matrix.py}; a completion audit and summary can
be regenerated with \path{rebuttal_experiments/aggregate.py}.  The separate
ten-seed E5 audit is \path{rebuttal_experiments/aggregate_e5n10.py}; it
requires the exact 380-file campaign before emitting any aggregate.

The original planned suite includes exactly 145 worker outputs.  E5
additionally reports the post-audit CUDA follow-ups under
\path{rebuttal_results/cuda/} and \path{rebuttal_results/cuda_tuning/}:
parallel tuning job 393048, phase profiles job 393047, production-path
stability job 393321, and RHS width/orientation probes jobs 393087 and 393161;
corrected CUDA test job 393425; inverse-factor
configuration-selection job 393427; and large-rank tuning job 393466.  Those
selection and diagnostic jobs contribute no final ten-seed latency.  All six
final E5 tables use campaign \texttt{e5n10-2ee46ba} from job 393591: one
33:48 preempt allocation on \texttt{n01}, ten serial seed blocks (1--10), and
a fresh Python process for every benchmark-driver invocation.  The refresh
driver emits five policy rows per invocation, and each large-rank pair
combines two method invocations.  The machine-readable campaign audit
reconciles all 380 tagged JSON files and every required row/seed key; a
separate source-to-\TeX{} audit reconciles every displayed E5 cell.  Earlier
measurement jobs 393470, 393428, 393454, 393472, 393491, and 393509 are
retained as superseded pilots but excluded from final ten-seed summaries.
Invalidated layout-bug jobs 393306 and 393301 are also explicitly excluded.
Files whose names contain \texttt{smoke}, tuning rows, and early reduced-rank
diagnostics remain excluded from inferential summaries and final speedup
claims.

The real-data ladder is outer campaign
\texttt{blackhole-s4-v1-locked-rerun2}, with raw results, attestations, audit,
and logs under
\path{rebuttal_results/blackhole/blackhole-s4-v1-locked-rerun2/}.
Its unmodified numerical driver and kernel execute from detached commit
\hexid{96c8018cdd0b56ffe1fd93bddf48fa4affb34697}; the rerun-2 control
tree is commit
\hexid{168478e0c02770937542bf045c7c13b36baa319f}, its locked harness
commit is
\hexid{653d8d916da04e1031c0ef2408093152a22642b4}, and the campaign-lock
SHA-256 is
\hexid{cecb0e4fe3271754dab72f86b43f4ee4ee53929ea35982661c99e376adf7f660}.
The raw HDF5, feature manifest, and feature-trace SHA-256 values are,
respectively,
\hexid{21e587c90c006c4b7d37e2b04899bd7e928967ae506e9ceb5afd05dfca6909fd},
\hexid{3b8c6e32da3842b8a9f21ee722e8b9fb1093899744d283d4e939427dd116a5c2},
and
\hexid{bdf1619a70af70585af317d25aa69452fcb8e94d8417e525de9d50f878e530b4}.

Slurm array 394033 ran under account \texttt{marlowe-m000151-pm06} in the
\texttt{batch} partition.  Its seven element job IDs used seven distinct
H100 UUIDs; every raw result records Python 3.10.12, PyTorch
2.12.0+\texttt{cu126}, CUDA 12.6, and cuDNN 91002.  Array task 3 records
\texttt{Restarts=1}: its first allocation failed during Slurm environment
retrieval before the batch script produced any output, claim, result, or
attestation; the same element then ran once scientifically on node
\texttt{n30}.  This infrastructure recovery is retained in provenance and
does not change task identity or the frozen protocol.  The original
mixed-HEAD diagnostic campaign is excluded.  Locked rerun~1 (array 393982)
is also excluded: its four-CPU lock rejected the scheduler's actual
12-CPU allocation before numerical execution, and its failed claims remain
preserved.  The terminal accounting and task-3 recovery record is
\path{rebuttal_results/blackhole/blackhole-s4-v1-locked-rerun2/SCHEDULER_PROVENANCE.md}.

The write-once aggregate was emitted only after all seven results and
attestations passed both the original v1 validator and the outer lock audit.
It treats the ten physical transitions as the statistical observations and
retains repeat-level samples without pseudoreplication.  The plotting script
\path{rebuttal_experiments/plot_blackhole_rank_scaling.py} accepts only this
passing campaign and reads
\path{blackhole-s4-v1-locked-rerun2/aggregate.json}.

Corrected E2 is a separate post-audit campaign with the following paths,
relative to the repository root:
\begin{quote}
\small
Root: \path{Cholesky-main-2/corrected_incremental_cholesky/}\\
Driver: \path{e2_budget_protein.py}\\
Slurm entry point: \path{slurm/e2_budget_protein.sbatch}\\
Artifacts: \path{results/e2_budget_protein/e2proper-fastpositive-v1/}
\end{quote}
Slurm job 393642 ran the three excluded pilots and twelve fixed evaluation
runs on \texttt{n02}; the campaign tag is
\texttt{e2proper-fastpositive-v1}.  All six \ours{} JSONs record source hash
\texttt{b1c92d7f...d87b4}.
The superseded \texttt{e2proper-dataf073752} campaign is retained only for
the explicitly confounded timing diagnostic and is not pooled with the final
evaluation; its raw JSONs survive, but its un-hashed historical executable
does not.  The recovered Protein data file is identified in the campaign
metadata by SHA-256 prefix \texttt{5ce538fa}; it is not a repository result
artifact.

\section{Limitations of the new evidence}

{\small
\begin{itemize}
  \item E1 uses a reconstructed natural-gradient backend and does not
  reproduce submitted Setup~B; it should not replace the submitted
  application results.
  \item Corrected E2 uses the recovered Protein data but a newly reconstructed
  driver rather than the missing submitted executable.  It has only three
  evaluation seeds, and the ten-epoch cap leaves most of the frozen time
  ceiling unused for \refact{} at \(r=256\) and \(r=512\), so only the
  matched-\(r=1024\) result supports the causal accuracy comparison.  The
  timing change from the prior diagnostic is also confounded by host
  (\texttt{n05} versus \texttt{n02}) and floating-point trajectory changes.
  \item E3 evaluates general fixed-width \(M=BC^\top\), not the symmetric
  specialization or a single online-growing-factor trajectory.
  \item E4 uses synthetic controlled spectra and RBF kernels.  It isolates
  numerical mechanisms but is not an application-level failure-rate study.
  \item E5's production-compatible batch-one comparison is limited to one
  H100, FP32, and \(r\le1024\).  The inverse-factor path changes stored state
  and is evaluated on ten synthetic problem seeds; its repeated-update and
  refresh validation stop at \(r=1024\), and both large-rank ladders measure
  isolated prepared updates.  All ten synthetic seeds ran serially on one
  GPU, so their SD does not measure device-to-device variation.  The real-data
  ladder uses one distinct H100 per rank and ten temporally related
  transitions from one trajectory, so its intervals are descriptive rather
  than estimates across independent simulations and its cross-rank curve
  includes device variation.  It excludes HDF5/DCT preprocessing,
  host-to-device transfer, and end-to-end forecasting, and does not include a
  specialized Woodbury baseline or the preregistered 100-transition secondary
  trace.  Periodic refresh cost is included only in the \(r\le1024\)
  synthetic refresh sweep.  Both ladders compare full dense results through
  \(r=32{,}768\); at larger ranks the synthetic ladder checks full-output
  finiteness, whereas the real-data ladder additionally compares sixteen
  fresh action probes.  Rank \(65{,}536\) requires 64\,GiB of matrix state and
  leaves limited device-memory headroom.
  \item E6 compares repeated rank-one calls with refactorization; it does not
  evaluate a native blocked rank-\(k\) kernel.
\end{itemize}
}

\end{document}
